\documentclass{article}
\usepackage[top=1in, bottom=1in, left=1in, right=1in]{geometry}
\usepackage{amsmath,amsthm,amssymb,amsfonts,mathtools,bm,bbm,array,graphicx,float,mathdots,dsfont,mathrsfs,latexsym}
\usepackage{caption,subcaption,multirow,longtable,lscape,wrapfig,tikz,tikz-cd,comment}
\captionsetup{font=footnotesize}
\usepackage[colorlinks=true, pdfstartview=FitV,linkcolor=blue,citecolor=blue,urlcolor=blue]{hyperref}
\usepackage[toc,page]{appendix}
\usepackage{booktabs}
\usepackage{multirow}
\usepackage{algorithm}
\usepackage{algpseudocode}
\DeclareMathOperator{\Z}{\mathbb{Z}}
\usepackage{blkarray}
\usepackage{amsmath}


\newtheorem{theorem}{Theorem}[section]
\newtheorem{conjecture}{Conjecture}[section]

\renewcommand{\thefootnote}{\fnsymbol{footnote}}

\newcommand{\hflogo}{\raisebox{-0.2em}{\includegraphics[height=1em]{figures/huggingface-color.png}}}
\newcommand{\ghlogo}{\raisebox{-0.2em}{\includegraphics[height=1em]{figures/github.png}}}


% --- theorem environments and macros for this article ---
\newtheorem{lemma}[theorem]{Lemma}
\newtheorem{proposition}[theorem]{Proposition}
\newtheorem{corollary}[theorem]{Corollary}
\newtheorem*{maintheorem}{Main Theorem}
\newtheorem*{theoremB}{Theorem B}
\newtheorem*{theoremC}{Theorem C}
\newtheorem*{corollaryD}{Corollary D}
\theoremstyle{definition}
\newtheorem{remark}[theorem]{Remark}
\newtheorem{definition}[theorem]{Definition}
\newtheorem{construction}[theorem]{Construction}
\theoremstyle{plain}
\newcommand{\Ric}{\operatorname{Ric}}
\newcommand{\R}{\mathbb{R}}
\newcommand{\HH}{\mathbb{H}}
\newcommand{\ip}[2]{\langle #1,#2\rangle}
\newcommand{\artanh}{\operatorname{artanh}}
\newcommand{\sech}{\operatorname{sech}}
\newcommand{\Vs}{\mathcal V_{\mathrm s}}
\newcommand{\Vb}{\mathcal V_{\mathrm b}}
\newcommand{\C}{\mathbb{C}}
\newcommand{\Q}{\mathbb{Q}}
\newcommand{\PP}{\mathbb{P}}
\newcommand{\cX}{\mathcal{X}}
\newcommand{\cO}{\mathcal{O}}
\newcommand{\eb}{\bar e}
\newcommand{\Ut}{\widetilde{U}}
\newcommand{\Yo}{Y^{\circ}}
\DeclareMathOperator{\Ima}{Im}
\DeclareMathOperator{\Rea}{Re}
\DeclareMathOperator{\Aut}{Aut}

\title{An Explicit Complex Structure on $S^4 \times T^2$}
\author{Edward Hirst\textsuperscript{1}\footnote{\href{mailto:ehirst@unicamp.br}{ehirst@unicamp.br}}, 
Tancredi Schettini Gherardini\textsuperscript{2,3}\footnote{\href{mailto:tsg@math.uni-bonn.de}{tsg@math.uni-bonn.de}}\\
Tom\'as S. R. Silva\textsuperscript{1}\footnote{\href{mailto:tomas@ime.unicamp.br}{tomas@ime.unicamp.br}},
Jakob Stein\textsuperscript{1}\footnote{\href{mailto:jstein@unicamp.br}{jstein@unicamp.br}}}
\date{\small\today\\ \mbox{}\\
\raggedright
\textit{\small \qquad\qquad \textsuperscript{1} Instituto de Matemática, Estatística e Computação Científica (IMECC), Universidade Estadual de \\\qquad\qquad\qquad\qquad\qquad\qquad Campinas (UNICAMP), 13083-859, Campinas, S\~ao Paulo, Brazil\\
\qquad\qquad\qquad\qquad\qquad \textsuperscript{2} Mathematisches Institut (MI), University of Bonn, Bonn, Germany\\
\qquad\qquad\qquad\qquad\qquad \textsuperscript{3} Max Planck Institute for Mathematics, Bonn, Germany \\
}}

\begin{document}
\maketitle
% \begingroup
% \renewcommand{\thefootnote}{}
% \footnotetext{Authors listed alphabetically.}
% \endgroup
\renewcommand{\thefootnote}{\arabic{footnote}}
\setcounter{footnote}{0}

\begin{abstract}
We construct a compact complex threefold $Y$ diffeomorphic to $S^4\times S^1\times S^1$.
Hence the six-manifold $S^4\times T^2$ admits a complex structure, which answers
affirmatively, for $V=T^2$, Calabi's question whether $V^2\times S^4$ admits a complex
structure. The threefold carries a holomorphic map $f\colon Y\to\PP^1$. Its fibres over
$\C^*$ are complex $2$-tori, and its two singular fibres are $E\times I_1$ (an elliptic
curve times a nodal rational curve). It also carries a free holomorphic $T^2$-action,
and its algebraic dimension is one.
The construction uses a holomorphic family of complex $2$-tori over $\C^*$ whose two
cusps have \emph{opposite chirality}: the elliptic curve at one cusp is oriented
oppositely to the one at the other. Such a family cannot be polarised, since polarised
families over $\C^*$ have no cusps, and its very general member is a torus with
trivial N\'eron--Severi group. The family is
completed by Mumford--Tate fillings and one Kodaira logarithmic transform of
multiplicity one. The quotient by the $T^2$-action is a simply connected homotopy
$4$-sphere, and the identification with $S^4\times T^2$ follows from Wall's
$h$-cobordism theorem and the $s$-cobordism theorem. By Iwase's classification of
torus fibrations, the quotient is in fact diffeomorphic to $S^4$.
We also run the Calabi--Eckmann construction over the complex $S^2\times S^4$ of Liang,
Liu and Zhang, using a holomorphic line bundle whose Chern class generates $H^2$. This
gives complex structures on $S^4\times S^3\times S^{2q+1}$ for every $q\ge0$. Combined
with the Main Theorem and the 2026 constructions on $S^6$ and $S^2\times S^4$, it follows
that every almost complex product of spheres of dimension at most eight is complex.


%\small Code: \ghlogo \, \href{...}{\texttt{...}}.
\end{abstract}
\newpage
\numberwithin{equation}{section}
\numberwithin{figure}{section}
\numberwithin{table}{section}
\tableofcontents
\newpage

%======================================================================
\section{Introduction}
%======================================================================

\subsection{Main result}
Calabi observed that $V^2\times S^4$ is almost complex for every closed orientable
surface $V^2$, and asked whether it admits a complex structure \cite{Calabi58}; see
\cite[\S1]{GFL26} and \cite{EuhSekigawa13} for recent accounts. The case $V=S^2$
was treated in 2026 by Guo, Fang and Lu \cite{GFL26} and by Liang, Liu and Zhang
\cite{LLZ26}. Both build on the 2026 complex structures on $S^6$
\cite{AC26,Engel26}, and \cite{LLZ26} explicitly assume the validity of the
construction of \cite{AC26}. The Main Theorem does not depend on any of these works:
its proof uses only the classical results listed in \S\ref{sec:facts} and the explicit
computations below.

\begin{maintheorem}
There is a compact complex manifold $Y$ of complex dimension $3$ that is diffeomorphic
to $S^4\times S^1\times S^1$. In particular, the smooth manifold $S^4\times T^2$
admits an integrable almost complex structure.
\end{maintheorem}

\begin{theoremB}
The threefold $Y$ of the Main Theorem has the following properties.
\begin{enumerate}
\item There is a proper surjective holomorphic map $f\colon Y\to\PP^1$. Its fibres
  over $\C^*$ are complex $2$-tori, and its fibres over $0$ and $\infty$ are
  isomorphic to $I_1\times E_0$ and $I_1\times E_\infty$. Here $I_1$ is a nodal
  rational curve and $E_0\cong E_\infty\cong\C/(\Z+i\Z)$. The two cusps differ in the
orientation of the marked lattice bases, $(1,i)$ and $(-1,i)$.
\item The real torus $T^2$ acts freely and holomorphically on $Y$, preserving $f$.
  The quotient $W=Y/T^2$ is a smooth homotopy $4$-sphere, and $Y\cong W\times T^2$
  as principal $T^2$-bundles. The complex Lie group $\C^2$ acts holomorphically on
  $Y$ by fibrewise translations, transitively on the smooth fibres of $f$.
\item $c_1(Y)=0$ in $H^2(Y;\Z)$, and $K_Y\cong f^*\cO_{\PP^1}(-4)$. In particular $K_Y$
  is topologically trivial but not a torsion element of $\mathrm{Pic}(Y)$.
\item The smooth fibres of $f$ are null-homologous in $H_4(Y;\Z)$. Consequently $Y$ is
  neither K\"ahler nor balanced.
\item Every meromorphic function on $Y$ is pulled back from $\PP^1$, i.e.\
  $\mathcal M(Y)=f^*\C(\PP^1)$. In particular the algebraic dimension of $Y$ is $1$.
\end{enumerate}
In addition, $W$ is diffeomorphic to the standard $S^4$
(Proposition~\ref{prop:Wstandard}); this uses Iwase's classification of torus
fibrations \cite{Iwase84}, and gives a second proof of the Main Theorem.
\end{theoremB}

\subsection{Idea of the construction}
Matsumoto showed that $S^4$ is the total space of a genus-one \emph{achiral}
Lefschetz fibration over $S^2$ \cite{Matsumoto82}. It has two singular fibres of types
$I_1^+$ and $I_1^-$, of \emph{opposite} signs, with the same vanishing cycle $a$
\cite[Thm.~4.2]{Matsumoto85}; see also \cite[Ex.~8.4.7]{GompfStipsicz99}. The fibration
has no section: a twisted gluing kills the dual cycle $b$. Our quotient $W$ will carry
a fibration with the same singular fibres, the same monodromy and the same twist
class. Such a fibration
cannot be holomorphic, because a family of elliptic curves only has positive nodes.

For complex $2$-tori the situation is different. Write the fibre near a cusp as
$X_t=\C^2/\langle e_1,e_2,e_3,e_4\rangle$, with vanishing cycle $e_1=(2\pi,0)$. The
intermediate quotient $\C^2/\langle e_1,e_2,e_3\rangle$ is an extension of the
elliptic curve $E_t=\C/\langle\eb_2,\eb_3\rangle$ (spanned by the second coordinates
of $e_2,e_3$) by $\C^*=\C/2\pi\Z$, and $X_t$ is its quotient by the infinite cyclic
group generated by $e_4$. The compact torus $X_t$ itself is not such an extension: the
kernel of a homomorphism of compact complex tori is compact. As $t$ tends to the cusp,
$X_t$ degenerates to a central fibre $I_1\times E$, whose smooth locus is
$\C^*\times E$. The orientation of $E$ relative to the complex orientation of the
fibre is an extra sign, not fixed by holomorphy
(Remark~\ref{rem:chirality}). We write down an explicit holomorphic family of
$2$-tori over $\C^*$ (\S\ref{sec:family}) with three properties:
\begin{itemize}
\item the elliptic curve is $\C/\langle 1,i\rangle$ at $t=0$ and
  $\C/\langle -1,i\rangle$ at $t=\infty$, so the two cusps have opposite
  chirality;
\item the full rank-$4$ lattice nevertheless stays nondegenerate for every
  $t\in\C^*$;
\item the vanishing cycle is the same at both ends.
\end{itemize}
The second property depends on a null conic in the period map. The family cannot be
polarised, since polarised families over $\C^*$ have no cusps
(Remark~\ref{rem:polarised}). Individual fibres can still be projective: $X_1$ is a
product of elliptic curves. But outside a countable set of parameters the fibres
have trivial N\'eron--Severi group, and hence no curves and no non-constant
meromorphic functions (Proposition~\ref{prop:NS}).

We then do three things:
\begin{enumerate}
\item fill both cusps with Mumford--Tate models (\S\ref{sec:filling});
\item perform one logarithmic transform of multiplicity one, whose class is the
  coinvariant $[e_4]$ (\S\ref{sec:twist}), obtaining a compact complex threefold $Y$;
\item observe that translations in two monodromy-invariant lattice directions give
  a free holomorphic $T^2$-action (\S\ref{sec:torus}).
\end{enumerate}
The quotient $W=Y/T^2$ inherits a genus-one achiral fibration
$f_W\colon W\to\PP^1$ with one node of each sign and a common vanishing cycle. These
are the singular-fibre data of Matsumoto's fibration of $S^4$, but we do not identify
$f_W$ with Matsumoto's map. Instead we prove directly that $W$ is simply connected
with $\chi(W)=2$ (\S\ref{sec:topology}), hence a homotopy $4$-sphere.
The Main Theorem then follows from the triviality of the $T^2$-bundle, Wall's
$h$-cobordism theorem and the $s$-cobordism theorem (\S\ref{sec:proof}).
Alternatively, $f_W$ is a good torus fibration in the sense of Matsumoto, and Iwase's
classification of such fibrations \cite{Iwase84} identifies $W$ with the standard
$S^4$ (Proposition~\ref{prop:Wstandard}).

\subsection{Relation to the constructions on $S^6$ and $S^2\times S^4$}
The toolkit is the same as in \cite{AC26,Engel26,GFL26}:
\begin{itemize}
\item an explicit period map of complex $2$-tori over a punctured $\PP^1$, whose very
  general member is non-algebraic;
\item Mumford-type toric fillings of the cusps;
\item Kodaira logarithmic transforms \cite{Kodaira63};
\item $\pi_1$ via van Kampen, and a classification or cobordism theorem for the
  diffeomorphism type.
\end{itemize}
The $(3,4,\infty)$ modular family used there does not seem suited to $S^4\times T^2$.
Let $Y\to\PP^1$ be a $2$-torus fibration with $Y\simeq S^4\times T^2$. Assume that
global sections of $R^1f_*\Q$ are detected on a smooth fibre. This holds, for instance,
when every singular fibre is of the type $I_1\times E$ used here; we have not checked
it for other degenerations. Then $b_1=2$ forces all local monodromies to be
transvections with a common vanishing direction (Remark~\ref{rem:necessity}). In
particular, no local monodromy can be of finite order other than the identity, so the
monodromies of order $3$ and $4$ of the $(3,4,\infty)$ family are excluded under this
assumption. Our family is in this sense the simplest possibility: two cusps,
monodromies $T^{\pm1}$.

\subsection{Products of spheres}\label{sec:introspheres}
Calabi's question is part of the broader problem of deciding which products of spheres
admit complex structures. The known results are as follows.
\begin{itemize}
\item Products of two odd-dimensional spheres are complex: these are the Hopf and
  Calabi--Eckmann manifolds \cite{CalabiEckmann53}.
\item Among products of even-dimensional spheres, the almost complex ones are exactly
  the products of copies of $S^2$, $S^6$ and $S^2\times S^4$ \cite{Sutherland65}. Since
  2026 all of these are complex \cite{AC26,Engel26,GFL26,LLZ26}.
\item Products containing both an $S^4$-factor and odd-dimensional factors are almost
  complex, being parallelisable \cite{Staples67}, but were not covered by these
  constructions. $S^4\times T^2$ is the smallest of them.
\end{itemize}

In \S\ref{sec:spheres} we combine the Main Theorem with these results. One more input is
needed: the complex $S^2\times S^4$ of Liang, Liu and Zhang \cite{LLZ26}, a compact complex
threefold $X_+$, carries a holomorphic line bundle $L_+$ whose Chern class generates
$H^2(X_+;\Z)$. Running the Calabi--Eckmann construction over $X_+\times\PP^q$, with the
line bundles $L_+$ and $\cO_{\PP^q}(-1)$, gives the following.

\begin{theoremC}
Assume the input \textup{(I3)} of \S\ref{sec:inputs}, i.e.\ the complex
$S^2\times S^4$ of \cite{LLZ26} with its line bundle $L_+$. Then:
\begin{enumerate}
\item for every $q\ge0$ there is a holomorphic elliptic-curve bundle over
  $X_+\times\PP^q$ whose total space is diffeomorphic to $S^4\times S^3\times S^{2q+1}$;
\item there is a holomorphic elliptic-curve bundle over $X_+\times X_+$ whose total space
  is diffeomorphic to $S^4\times S^3\times S^4\times S^3$.
\end{enumerate}
In particular $S^4\times S^3\times S^1$ admits a complex structure.
\end{theoremC}

\begin{corollaryD}
Assume the inputs \textup{(I1)--(I3)} of \S\ref{sec:inputs}. Then a product of spheres
of dimension at most eight admits a complex structure if and only if it admits an
almost complex structure. The only such products with no almost complex structure are
$S^4$, $S^8$ and $S^4\times S^4$.
\end{corollaryD}

Corollary~D uses the Main Theorem for exactly two products: $S^4\times T^2$, and
$S^4\times T^4\cong Y\times E$ with $E$ an elliptic curve. It uses Theorem~C for
$S^4\times S^3\times S^1$. Without the Main Theorem, $S^4\times T^2$ would be the only open
case in dimension at most six, and $S^4\times T^4$ the only open case in dimension eight
(Remark~\ref{rem:wasopen}).

\subsection{Novelty}\label{sec:novelty}
As far as we have been able to determine (\S\ref{sec:literature}), no complex structure
on $S^4\times T^2$, or on $V^2\times S^4$ with $V$ of positive genus, was previously
known. The question for $V=T^2$ is recorded as open in \cite{Calabi58} and
\cite{MO89798}. Theorem~C is a direct combination of \cite{LLZ26} with the
Calabi--Eckmann construction. We have not seen it stated.

\subsection*{Conventions}
Manifolds are smooth, Hausdorff and second countable. We identify $\C^2$ with $\R^4$
via
\[(z,w)\mapsto(\Rea z,\Ima z,\Rea w,\Ima w);\]
this is the complex orientation, and $\det_\R$ denotes the real $4\times4$ determinant. For $u,v\in\C$ put
$\omega(u,v)=\Ima(\bar u v)=\Rea u\,\Ima v-\Ima u\,\Rea v$. An action of a discrete
group $\Gamma$ on a locally compact Hausdorff space $X$ is \emph{properly discontinuous}
if $\{\gamma:\gamma K\cap K\neq\emptyset\}$ is finite for every compact $K\subset X$. A free,
properly discontinuous action by biholomorphisms of a complex manifold has a complex
manifold as quotient, and the quotient map is a covering. $D_\varepsilon=\{|t|<\varepsilon\}$ and
$D^*_\varepsilon=D_\varepsilon\setminus\{0\}$.

%======================================================================
\section{Preliminaries}
%======================================================================

\subsection{Lattices in $\C^2$}

\begin{lemma}[determinant formula]\label{lem:det}
Let $e_1=(2\pi,0)$ and $e_j=(x_j,\eb_j)\in\C^2$ for $j=2,3,4$. Then
\[
\det\nolimits_\R(e_1,e_2,e_3,e_4)=2\pi\,\det\begin{pmatrix}\Ima x_2&\Ima x_3&\Ima x_4\\
\Rea\eb_2&\Rea\eb_3&\Rea\eb_4\\ \Ima\eb_2&\Ima\eb_3&\Ima\eb_4\end{pmatrix}
=2\pi\,\Ima X\cdot(\Rea\eb\times\Ima\eb),
\]
where $X=(x_2,x_3,x_4)$, $\eb=(\eb_2,\eb_3,\eb_4)$ and $\times$ is the cross product
in $\R^3$.
\end{lemma}
\begin{proof}
The real coordinate vector of $e_1$ is $(2\pi,0,0,0)$. Expand along the first column.
\end{proof}

\begin{lemma}[local determinant identity]\label{lem:localdet}
Let $\alpha_2,\alpha_3,\eb_2,\eb_3,\eb_4\in\C$ with $\omega(\eb_2,\eb_3)\ne0$,
and let $\ell\in\C$. Let $\varphi\colon\C\to\R$ be the $\R$-linear functional with
$\varphi(\eb_j)=\Rea\alpha_j$ for $j=2,3$, and put $P=\Rea\ell-\varphi(\eb_4)$. Then
\[
\det\nolimits_\R\big((2\pi,0),(-i\alpha_2,\eb_2),(-i\alpha_3,\eb_3),(-i\ell,\eb_4)\big)=
-2\pi\,P\,\omega(\eb_2,\eb_3).
\]
\end{lemma}
\begin{proof}
By Lemma~\ref{lem:det} the left side is $2\pi L(\Ima X)$, where
\[L(v)=\det[v;\Rea\eb;\Ima\eb],\qquad \Ima X=-(\Rea\alpha_2,\Rea\alpha_3,\Rea\ell).\]
\begin{itemize}
\item Write $\varphi(x+iy)=px+qy$. Then
  $u:=(\varphi(\eb_2),\varphi(\eb_3),\varphi(\eb_4))=p\Rea\eb+q\Ima\eb$, so $L(u)=0$.
\item Hence $L(\Ima X)=-L(u+(0,0,P))=-P\,L(0,0,1)$.
\item Finally $L(0,0,1)=\Rea\eb_2\Ima\eb_3-\Ima\eb_2\Rea\eb_3=\omega(\eb_2,\eb_3)$.
\end{itemize}
(This identity is also checked symbolically; see Appendix~\ref{app:checks}.)
\end{proof}

\begin{lemma}[families of lattices]\label{lem:latticefamily}
Let $B$ be a complex manifold, and let $v_1,\dots,v_4\colon B\to\C^2$ be holomorphic with
$\det_\R(v_1(b),\dots,v_4(b))\neq0$ for all $b$. Then $\Z^4$, acting on $\C^2\times B$ by
$(v,b)\mapsto(v+\sum n_jv_j(b),b)$, acts freely and properly discontinuously. The
quotient is a complex manifold, and the projection to $B$ is a proper holomorphic
submersion with fibres the complex tori $\C^2/\bigoplus_j\Z v_j(b)$.
\end{lemma}
\begin{proof}
Let $V(b)\in GL_4(\R)$ have columns $v_j(b)$.
\begin{itemize}
\item On a compact $L\subset B$, $\|V(b)^{-1}\|\le C_L$, so
  $|\sum n_jv_j(b)|\ge |n|/C_L$. Hence only finitely many $n$ move a compact subset of
  $\C^2\times L$ into itself. Freeness is clear.
\item The preimage of a compact $L$ is the image of the compact set
  $\{(\sum s_jv_j(b),b): s\in[0,1]^4,\ b\in L\}$, so the projection is proper.
\end{itemize}
\end{proof}

\subsection{The toric surface $\cX$}\label{sec:toric}
Let $N=\Z^2$, $v_n=(n,1)$ and $\sigma_n=\mathrm{cone}(v_n,v_{n+1})$ for $n\in\Z$. The
cones $\sigma_n$ and their faces form an infinite fan $\Sigma$, whose support is the
open upper half-plane together with the origin. Any two cones of $\Sigma$ meet in a
common face, and every compact subset of the open upper half-plane meets only finitely
many cones; this is the finiteness used in Mumford's construction \cite{Mumford72}. The
fan is not locally finite at the origin, which lies in every cone. Each
$\sigma_n$ is unimodular, since $\det(v_n,v_{n+1})=-1$. Let $\cX=\cX_\Sigma$ be the
associated smooth toric surface, obtained by gluing the affine charts
$U_{\sigma_n}\cong\C^2$ along the open subsets corresponding to common faces
\cite{Fulton93}. Since infinitely many charts are involved, we verify directly in
Lemma~\ref{lem:toric}(0) that the result is a Hausdorff, second countable complex
surface. Write $(\zeta,t)$ for the coordinates on its dense torus
$(\C^*)^2$, given by the characters $(1,0)$ and $(0,1)$.

\begin{lemma}\label{lem:toric}\leavevmode
\begin{enumerate}
\setcounter{enumi}{-1}
\item \emph{Separation.} $\cX$ is a Hausdorff, second countable complex manifold of
  dimension $2$.
\item \emph{Charts.} $\cX=\bigcup_n U_n$ with $U_n\cong\C^2$, coordinates
  $x_n=\zeta^{-1}t^{n+1}$ and $y_n=\zeta t^{-n}$. On $U_n$ we have $t=x_ny_n$ and
  $\zeta=x_n^ny_n^{n+1}$. On $U_n\cap U_{n+1}=\{x_n\neq0\}$ we have $x_{n+1}=tx_n$ and
  $y_{n+1}=1/x_n$.
\item \emph{The map $t$.} The character $t$ extends to a holomorphic function
  $t\colon\cX\to\C$. Every point with $t\neq0$ lies in the dense torus, so
  $t^{-1}(t_0)\cong\C^*_\zeta$ for $t_0\ne0$.
\item \emph{Central fibre.} $t^{-1}(0)=\bigcup_n D_n$, where $D_n\cong\PP^1$ is the
  divisor of $v_n$; $D_n\cap U_n=\{x_n=0\}$ and $D_{n+1}\cap U_n=\{y_n=0\}$. Thus
  $t^{-1}(0)$ is an infinite chain of $\PP^1$'s, with $D_n\cap D_{n+1}=\{P_n\}$ the
  origin of $U_n$, and $D_n\cap D_m=\emptyset$ for $|n-m|\ge2$.
\item \emph{Torus action.} $\lambda\in\C^*$ acts by $\lambda\cdot(\zeta,t)=(\lambda\zeta,t)$,
  i.e.\ $\lambda\cdot(x_n,y_n)=(\lambda^{-1}x_n,\lambda y_n)$. It preserves $t$, each $D_n$
  and each $P_n$.
\item \emph{Shift.} The lattice automorphism $A(a,b)=(a+b,b)$ maps $\sigma_n$ to
  $\sigma_{n+1}$. The induced automorphism $\psi$ of $\cX$ is $(\zeta,t)\mapsto(t\zeta,t)$
  on the torus; it maps $U_n$ to $U_{n+1}$ with $x_{n+1}\circ\psi=x_n$ and
  $y_{n+1}\circ\psi=y_n$. It commutes with the $\C^*$-action, preserves $t$, and maps
  $D_n$ to $D_{n+1}$ and $P_n$ to $P_{n+1}$.
\item \emph{Log form.} On $U_n$, $\tfrac{d\zeta}{\zeta}\wedge dt=-dx_n\wedge dy_n$.
\end{enumerate}
\end{lemma}
\begin{proof}
Items (1)--(6) are standard toric computations \cite{Fulton93}; the formulas in (1),
(5) and (6) are checked in Appendix~\ref{app:checks}.
\begin{itemize}
\item (1): the dual basis of $(v_n,v_{n+1})$ is
  $(m^{(n)}_1,m^{(n)}_2)=((-1,n+1),(1,-n))$, and $x_n=\chi^{m^{(n)}_1}$,
  $y_n=\chi^{m^{(n)}_2}$.
\item (2): $\langle(0,1),v_n\rangle=1\ge0$ for all $n$.
\item (5): with $A=\begin{pmatrix}1&1\\0&1\end{pmatrix}$ one has
  $A^Tm^{(n+1)}_1=m^{(n)}_1$ and $A^Tm^{(n+1)}_2=m^{(n)}_2$. Since
  $\chi^m\circ\psi=\chi^{A^Tm}$, this gives $x_{n+1}\circ\psi=x_n$ and
  $y_{n+1}\circ\psi=y_n$.
\end{itemize}
(0) Let $m=n+k$ with $k\ge1$, and write $t_j=x_jy_j$ on $U_j$. On the dense torus,
(1) gives $x_m=t^kx_n$ and $y_m=t^{-k}y_n$. Hence the set of pairs
$(p,q)\in U_n\times U_m$ that are identified in $\cX$ is
\[R_{n,m}=\big\{(p,q):\ t_n(p)=t_m(q),\ \ x_m(q)=t_n(p)^k\,x_n(p),\ \ t_n(p)^{k-1}x_n(p)\,y_m(q)=1\big\}.\]
\begin{itemize}
\item For $k=1$ this is the graph of the gluing map $(x_n,y_n)\mapsto(tx_n,1/x_n)$ on
  the overlap $\{x_n\neq0\}$.
\item For $k\ge2$ the last equation forces $t\ne0$ and $x_n\ne0$, so $R_{n,m}$ is the
  graph of the identification of the two dense tori, which is the overlap
  $U_n\cap U_m$ (the cones meet only in $\{0\}$).
\end{itemize}
Each $R_{n,m}$ is defined by polynomial equations, so it is closed in $U_n\times U_m$.
By symmetry the same holds for $m<n$, and $R_{n,n}$ is the diagonal. Thus the
equivalence relation on $\bigsqcup_nU_n$ is closed. The quotient map is open,
since the identifications are by homeomorphisms between open subsets. A quotient of a
Hausdorff space by an open map with closed equivalence relation is Hausdorff. The atlas
$\{U_n\}_{n\in\Z}$ is countable, so $\cX$ is second countable.
\end{proof}

For $0<\varepsilon<1$ put $\cX^\varepsilon=t^{-1}(D_\varepsilon)$, and for $\rho\ge1$ define the
\emph{slabs}
\[W_n(\rho)=\{p\in U_n:\ |x_n(p)|<\rho,\ |y_n(p)|<\rho,\ |t(p)|<\varepsilon\}.\]

\begin{lemma}[slabs]\label{lem:slabs}
Let $0<\varepsilon<1/16$.
\begin{enumerate}
\item The slabs $W_n(2)$ cover $\cX^\varepsilon$.
\item If $|n-m|\ge2$, then $W_n(2)\cap W_m(2)=\emptyset$.
\item Each $W_n(2)$ is contractible. Each $W_n(2)\cap W_{n+1}(2)$ is connected and
  homotopy equivalent to $S^1$.
\item $\psi(W_n(\rho))=W_{n+1}(\rho)$, and $\lambda\cdot W_n(\rho)\subset W_n(\rho\max(|\lambda|,|\lambda|^{-1}))$.
\item If $p\in W_n(2)$ and $t(p)\neq0$, then
  $\frac{\log|\zeta(p)|}{\log|t(p)|}\in(n-\tfrac12,\,n+\tfrac32)$.
\end{enumerate}
\end{lemma}
\begin{proof}
(5) On the torus, $|x_n|<2$ means $|\zeta|>|t|^{n+1}/2$, and $|y_n|<2$ means
$|\zeta|<2|t|^n$. Take logarithms, divide by $\log|t|<0$, and use
$\log2/|\log|t||<\tfrac12$.

(1) Let $p$ be a torus point. Put $n=\lfloor\log|\zeta|/\log|t|\rfloor$. Then
$|t|^{n+1}<|\zeta|\le|t|^n$, so $|x_n|<1$ and $|y_n|\le1$.
Now let $p\in D_n$, so $p\in U_{n-1}\cup U_n$.
\begin{itemize}
\item If $p\in U_n$ and $|y_n(p)|<2$, then $p\in W_n(2)$.
\item Otherwise either $p=P_{n-1}$, or $p\in U_n$ with $|y_n(p)|\ge2$. In the second
  case $p\in U_{n-1}\cap U_n$ and $|x_{n-1}(p)|=1/|y_n(p)|\le\tfrac12$. In both cases
  $p\in U_{n-1}$ with $y_{n-1}(p)=0$ and $|x_{n-1}(p)|\le\tfrac12$, so $p\in W_{n-1}(2)$.
\end{itemize}

(2) By (3) of Lemma~\ref{lem:toric}, $W_n(2)\cap t^{-1}(0)\subset D_n\cup D_{n+1}$, and
$P_m\in W_n$ only for $m=n$. So the central parts are disjoint for $|n-m|\ge2$.
For torus points with $m\ge n+2$, a common point would need
$|t|^{n+1}/2<|\zeta|<2|t|^m$, i.e.\ $|t|^{m-n-1}>1/4$. This is impossible since
$|t|<1/16$.

(3) In the coordinates $(x,y)=(x_n,y_n)$, $W_n(2)=\{|x|,|y|<2,\ |xy|<\varepsilon\}$ is
star-shaped about $0$. By Lemma~\ref{lem:toric}(1),
\[W_n(2)\cap W_{n+1}(2)=\{\tfrac12<|x|<2,\ |y|<\min(2,\varepsilon/|x|)\},\]
which retracts onto $\{|x|=1,y=0\}$.

(4) This follows from Lemma~\ref{lem:toric}(4),(5).
\end{proof}

\begin{lemma}[topology of $\cX^\varepsilon$]\label{lem:Xtop}
For $0<\varepsilon<1/16$, $\cX^\varepsilon$ is connected and simply connected, and
\[H_q(\cX^\varepsilon;\Z)\cong\Z,\ 0,\ \Z[\tau,\tau^{-1}],\ 0\quad\text{for } q=0,\ 1,\ 2,\ \ge3,\]
where $\tau$ acts on $H_2$ through $\psi_*$. Moreover, any homeomorphism
$p\mapsto\kappa(t(p))\cdot\psi(p)$, with $\kappa\colon D_\varepsilon\to\C^*$ admitting a continuous logarithm, acts on homology as $\psi_*$.
\end{lemma}
\begin{proof}
Put $Z_N=\bigcup_{|n|\le N}W_n(2)$. By Lemma~\ref{lem:slabs}, adding $W_{N+1}(2)$ to
$Z_N$ (or $W_{-N-1}(2)$) attaches a contractible set along a connected set
$\simeq S^1$.
\begin{itemize}
\item Van Kampen shows that each $Z_N$ is simply connected.
\item Mayer--Vietoris shows that each such step adds one free generator to $H_2$
  and changes nothing else. Explicitly, when $W_{n+1}$ is attached to $Z_n$ (along
  $W_n\cap W_{n+1}$), the boundary map $H_2(Z_n\cup W_{n+1})\to H_1(W_n\cap W_{n+1})\cong\Z$ is
  onto, and $H_2(Z_n)$ is its kernel.
\item The sphere $D_{n+1}$ lies in $W_n\cup W_{n+1}$: on $D_{n+1}$, $y_{n+1}=1/x_n$, so
  every point has $|x_n|<2$ or $|y_{n+1}|<2$. The circle $\{|x_n|=1\}\subset D_{n+1}$
  cuts it into two discs, one in each slab. So the boundary map sends $[D_{n+1}]$ to a
  generator of $H_1(W_n\cap W_{n+1})$.
\end{itemize}
Hence $H_2(Z_N)$ is free on the classes $[D_m]$ of the compact components in
$Z_N$, and $H_2(\cX^\varepsilon)$, the direct limit, is free on $\{[D_n]\}_{n\in\Z}$.
Since $\psi$ is holomorphic and maps $D_n$ to $D_{n+1}$, it preserves complex
orientations, so $\psi_*[D_n]=[D_{n+1}]$. Thus $H_2\cong\Z[\tau^{\pm1}]$ as a
$\Z[\tau^{\pm1}]$-module, with $[D_n]\leftrightarrow\tau^n$.
The map $p\mapsto\kappa(t(p))^u\cdot p$, with $\kappa^u=\exp(u\log\kappa)$ and
$u\in[0,1]$, is an isotopy from the identity to $p\mapsto\kappa(t(p))\cdot p$, which
gives the last statement.
\end{proof}

\subsection{Standard facts used}\label{sec:facts}
\begin{enumerate}
\item[(F1)] Principal $T^k$-bundles over a CW complex $B$ are classified by
  $H^2(B;\Z^k)$, and topologically trivial smooth principal bundles are smoothly
  trivial.
\item[(F2)] (Milnor \cite{Milnor68}) For an infinite cyclic covering $\widetilde X\to X$
  with deck generator $\tau$ there is an exact sequence
  $\cdots\to H_q(\widetilde X)\xrightarrow{\tau_*-1}H_q(\widetilde X)\to H_q(X)\to
  H_{q-1}(\widetilde X)\to\cdots$.
\item[(F3)] (Wall \cite{Wall64}) Two smooth closed simply connected $4$-manifolds with
  isomorphic intersection forms are $h$-cobordant. For homotopy spheres this is the
  statement $\Theta_4=0$ \cite{KervaireMilnor63}, where $\Theta_4$ is the group of
  $h$-cobordism classes of homotopy $4$-spheres. It does not assert the smooth
  $4$-dimensional Poincar\'e conjecture.
\item[(F4)] The Whitehead group $\mathrm{Wh}(\Z^2)$ vanishes \cite{BHS64}. The
  $s$-cobordism theorem \cite{Kervaire65,Milnor66} says an $h$-cobordism of dimension $\ge6$
  with vanishing Whitehead torsion is a product.
\end{enumerate}

%======================================================================
\section{The family of $2$-tori over $\C^*$}\label{sec:family}
%======================================================================

\begin{construction}\label{con:family}
For $t\in\C^*$ consider the following holomorphic automorphisms of
$\C^*_\zeta\times\C_w\times\C^*_t$, all of which preserve $t$:
\begin{align*}
g_2(\zeta,w,t)&=(e^{t}\zeta,\ w+1-t^2,\ t),\\
g_3(\zeta,w,t)&=(e^{-it}\zeta,\ w+i(1+t^2),\ t),\\
g_4(\zeta,w,t)&=(t\zeta,\ w+2t,\ t).
\end{align*}
They commute. Let $\Gamma^\circ=\langle g_2,g_3,g_4\rangle$ and $\Yo=(\C^*\times\C\times\C^*)/\Gamma^\circ$.
\end{construction}

In the additive coordinate $z$, $\zeta=e^{iz}$, and for a choice of $\log t$, the
lifts of $g_2,g_3,g_4$, together with $z\mapsto z+2\pi$, are the translations by
\begin{equation}\label{eq:lattice}
e_1=(2\pi,0),\quad e_2(t)=(-it,\,1-t^2),\quad e_3(t)=(-t,\,i(1+t^2)),\quad e_4(t)=(-i\log t,\,2t).
\end{equation}
Changing the branch of $\log t$ by $2\pi i$ replaces $e_4$ by $e_4+e_1$, so
$\Lambda_t=\bigoplus_{j}\Z e_j(t)\subset\C^2$ is independent of the branch.

\begin{lemma}\label{lem:det-family}
For all $t\in\C^*$, with $r=|t|$,
\[
\det\nolimits_\R(e_1,e_2,e_3,e_4)=2\pi(1+r^2)\big(2r^2-(1-r^2)\ln r\big)>0.
\]
\end{lemma}
\begin{proof}
Write $t=a+ib$. We have $\Ima x_2=-a$, $\Ima x_3=-b$, $\Ima x_4=-\ln r$, and
\[
\Rea\eb=(1-a^2+b^2,\,-2ab,\,2a),\qquad \Ima\eb=(-2ab,\,1+a^2-b^2,\,2b).
\]
A direct computation gives
$\Rea\eb\times\Ima\eb=(1+r^2)\cdot(-2a,\,-2b,\,1-r^2)$, so by Lemma~\ref{lem:det}
\[
\det\nolimits_\R=2\pi(1+r^2)(2a^2+2b^2-(1-r^2)\ln r).
\]
The factors $1-r^2$ and $\ln r$ have opposite signs for $r\ne1$ and vanish only at
$r=1$, so $(1-r^2)\ln r\le0$. Hence $\det_\R\ge 4\pi r^2(1+r^2)>0$.
\end{proof}

\begin{lemma}\label{lem:Yo}
$\Gamma^\circ\cong\Z^3$ acts freely and properly discontinuously, and $\Yo$ is a complex
manifold. The map $f\colon\Yo\to\C^*$, $[\zeta,w,t]\mapsto t$, is a proper holomorphic
submersion with fibres $X_t=\C^2/\Lambda_t$, which are compact complex $2$-tori. The
monodromy of $\{H_1(X_t;\Z)=\Lambda_t\}$ around $t=0$ (counterclockwise) is
\[
T\colon e_1\mapsto e_1,\quad e_2\mapsto e_2,\quad e_3\mapsto e_3,\quad e_4\mapsto e_4+e_1.
\]
The zero section $t\mapsto[1,0,t]$ is a holomorphic section of $f$.
\end{lemma}
\begin{proof}
Let $B\subset\C^*$ be a simply connected open disc and fix $\log t$ on $B$. The map
$\C^2\times B\to\C^*\times\C\times B$, $(z,w,t)\mapsto(e^{iz},w,t)$, is the quotient
by translation by $e_1$. Under it, $\Gamma^\circ$ corresponds to the translations by
$e_2,e_3,e_4$ modulo $e_1$.
\begin{itemize}
\item By Lemmas~\ref{lem:latticefamily} and \ref{lem:det-family}, the $\Z^4$-action
  on $\C^2\times B$ is free and properly discontinuous, with quotient
  $f^{-1}(B)\to B$ a proper holomorphic submersion with fibres $\C^2/\Lambda_t$.
\item Since $\Gamma^\circ$ preserves $t$, proper discontinuity and freeness on all of
  $\C^*\times\C\times\C^*$ follow by covering a compact set's image in $\C^*$ by
  finitely many such discs.
\item $\Gamma^\circ\cong\Z^3$ because $e_1,\dots,e_4$ are $\R$-independent.
\item The monodromy statement is the branch change above.
\end{itemize}
\end{proof}

\begin{remark}[opposite chirality]\label{rem:chirality}
Call the second coordinates $\eb_2(t)=1-t^2$ and $\eb_3(t)=i(1+t^2)$ of $e_2,e_3$ the
two \emph{projected period vectors}. They satisfy
$\omega(\eb_2,\eb_3)=1-|t|^4$.
\begin{itemize}
\item For $|t|\ne1$ they span a lattice in $\C$, with modulus
  $\tau_E(t)=\eb_3/\eb_2=i(1+t^2)/(1-t^2)$. This is $i$ at $t=0$ and $-i$ at
  $t=\infty$. Near the two cusps these are the lattices of the elliptic curves $E_0$
  and $E_\infty$, which therefore have opposite orientations (Lemma~\ref{lem:atinfty}).
\item On $|t|=1$ they are $\R$-linearly dependent, so they do not span a lattice. For
  $t=e^{i\theta}$ with $\theta\notin\pi\Z$, $\tau_E=-\cot\theta$ is real, and if
  $\cot\theta$ is irrational they span a dense subgroup of a real line. At $t=\pm1$,
  $\eb_2=0$.
\end{itemize}
By Lemma~\ref{lem:det-family}, the full rank-$4$ lattice $\Lambda_t$ nevertheless stays
nondegenerate across $|t|=1$.

The vector $\eb(t)=(1-t^2,\,i(1+t^2),\,2t)$ parametrises the null conic
$\{x^2+y^2+z^2=0\}\subset\PP^2$. This conic has no real points, which is what makes
$\Rea\eb\times\Ima\eb$ nowhere zero.
\end{remark}

\subsection{Periods and N\'eron--Severi groups}
Let $B\subset\C^*$ be an open disc, and fix a branch of $\log t$ on $B$, so that
$e_1,\dots,e_4$ is a basis of $\Lambda_t=H_1(X_t;\Z)$ depending holomorphically on
$t\in B$. For $j<k$ let $\Sigma_{jk}\subset X_t$ be the $2$-torus spanned by $e_j$ and
$e_k$, oriented by $(e_j,e_k)$. The classes $[\Sigma_{jk}]$ form a basis of
$H_2(X_t;\Z)=\Lambda^2\Lambda_t$. The \emph{periods} of the holomorphic $2$-form
$\sigma=dz\wedge dw$ are
\[p_{jk}(t)=\int_{\Sigma_{jk}}\sigma=z_jw_k-z_kw_j,\qquad e_j=(z_j,w_j).\]

\begin{lemma}[periods]\label{lem:periods}
\begin{enumerate}
\item The periods are
\begin{align*}
p_{12}&=2\pi(1-t^2), & p_{13}&=2\pi i(1+t^2), & p_{14}&=4\pi t,\\
p_{23}&=2t, & p_{24}&=i(1-t^2)\log t-2it^2, & p_{34}&=-(1+t^2)\log t-2t^2.
\end{align*}
\item If $c=(c_{jk})\in\Z^6$ and $\sum_{j<k}c_{jk}p_{jk}$ vanishes identically on $B$,
  then $c=0$.
\end{enumerate}
\end{lemma}
\begin{proof}
(1) Parametrise $\Sigma_{jk}$ by $(a,b)\mapsto ae_j+be_k$, $(a,b)\in[0,1]^2$. Then
$dz\wedge dw$ pulls back to $(z_jw_k-z_kw_j)\,da\wedge db$. Insert \eqref{eq:lattice};
see also Appendix~\ref{app:checks}.

(2) Write $F=\sum c_{jk}p_{jk}=A(t)+C(t)\log t$ with polynomials
\begin{align*}
C(t)&=ic_{24}(1-t^2)-c_{34}(1+t^2),\\
A(t)&=2\pi c_{12}(1-t^2)+2\pi ic_{13}(1+t^2)+(4\pi c_{14}+2c_{23})\,t-2(ic_{24}+c_{34})\,t^2.
\end{align*}
\begin{itemize}
\item If $C\not\equiv0$, then $\log t=-A/C$ on $B$. Continuing analytically once around
  $0$ changes the left side by $2\pi i$ and leaves the right side unchanged, which is
  impossible. So $C\equiv0$, and then $A\equiv0$.
\item The constant term of $C$ is $ic_{24}-c_{34}$. Since $c_{24},c_{34}$ are real,
  $c_{24}=c_{34}=0$.
\item The constant term of $A$ is $2\pi(c_{12}+ic_{13})$, so $c_{12}=c_{13}=0$.
\item The linear term of $A$ is $4\pi c_{14}+2c_{23}$. Since $\pi$ is irrational,
  $c_{14}=c_{23}=0$.
\end{itemize}
\end{proof}

\begin{proposition}[N\'eron--Severi groups of the fibres]\label{prop:NS}
\begin{enumerate}
\item There is a countable subset $\mathcal E$ of $\C^*$ with the following property.
  For every $t\in\C^*\setminus\mathcal E$ the N\'eron--Severi group of $X_t$ is zero,
  $X_t$ contains no curves, and $X_t$ has no non-constant meromorphic functions, i.e.\
  $a(X_t)=0$.
\item The fibre over $t=1$ is a product of elliptic curves,
  \[X_1\cong\C/(2\pi\Z+i\Z)\times\C/(2\Z+2i\Z),\]
  hence a projective abelian surface. In particular $1\in\mathcal E$, so not every
  smooth fibre is non-algebraic.
\end{enumerate}
\end{proposition}
\begin{proof}
(1) By the Lefschetz $(1,1)$-theorem, $\mathrm{NS}(X_t)$ is the group of integral
classes of type $(1,1)$.
\begin{itemize}
\item A real class $\gamma\in H^2(X_t;\R)$ is of type $(1,1)$ if and only if
  $\int_{X_t}\gamma\wedge\sigma=0$. Indeed, only the $(0,2)$-part of $\gamma$
  contributes to this integral. That part is a multiple of $\bar\sigma$, and
  $\int\bar\sigma\wedge\sigma\ne0$. Moreover, $\gamma^{2,0}=\overline{\gamma^{0,2}}$.
\item Write the Poincar\'e dual of an integral class $\gamma$ as
  $\sum c_{jk}[\Sigma_{jk}]$ with $c\in\Z^6$, using the convention
  $\int_{X_t}\alpha\wedge\gamma=\int_{\mathrm{PD}(\gamma)}\alpha$. Then
  $\int\gamma\wedge\sigma=\sum c_{jk}p_{jk}(t)$.
\item For $c\ne0$, this holomorphic function on $B$ is not identically zero by
  Lemma~\ref{lem:periods}. So its zeros form a discrete, hence countable, subset of
  $B$.
\item $\C^*$ is covered by countably many discs $B$, and $\Z^6$ is countable. So the
  set $\mathcal E$ of $t$ with $\mathrm{NS}(X_t)\ne0$ is countable.
\end{itemize}
Now let $t\notin\mathcal E$. A curve $C\subset X_t$ would have $[C]\ne0$ in
$\mathrm{NS}(X_t)$, since $\int_C\omega_0>0$ for a flat K\"ahler form $\omega_0$ on
$X_t$. A non-constant meromorphic function has a non-empty polar divisor. So neither
exists.

(2) At $t=1$ the generators are $e_1=(2\pi,0)$, $e_2=(-i,0)$, $e_3=(-1,2i)$ and
$e_4=(0,2)$. In the coordinates $(z',w)=(z-\tfrac i2w,\ w)$ they become $(2\pi,0)$,
$(-i,0)$, $(0,2i)$ and $(-i,2)$. Replacing the last vector by $e_4-e_2=(0,2)$ exhibits
$\Lambda_1$ as the product lattice $(2\pi\Z+i\Z)\times(2\Z+2i\Z)$.
\end{proof}

%======================================================================
\section{Filling the cusps}\label{sec:filling}
%======================================================================

\begin{proposition}[Mumford--Tate filling]\label{prop:filling}
Let $\varepsilon_0\in(0,1/16)$, let $\alpha_2,\alpha_3,\eb_2,\eb_3,\eb_4$ be holomorphic on
$D_{\varepsilon_0}$, and let $\nu\in\C^*$, such that
\begin{enumerate}
\item[(H1)] $\alpha_2(0)=\alpha_3(0)=\eb_4(0)=0$;
\item[(H2)] $\omega(\eb_2(t),\eb_3(t))\neq0$ for $t\in D_{\varepsilon_0}$.
\end{enumerate}
On $\cX^{\varepsilon_0}\times\C_w$, where $\cX^{\varepsilon_0}=t^{-1}(D_{\varepsilon_0})$
and $t=t(p)$, define
\[
G_j(p,w)=(e^{\alpha_j(t)}\cdot p,\ w+\eb_j(t))\ \ (j=2,3),\qquad
G_4(p,w)=(\nu\cdot\psi(p),\ w+\eb_4(t)).
\]
There is $\varepsilon\in(0,\varepsilon_0)$ with the following properties, where
$\Ut=\cX^\varepsilon\times\C$ and $\Gamma=\langle G_2,G_3,G_4\rangle$.
\begin{enumerate}
\item[(a)] $\Gamma\cong\Z^3$ acts on $\Ut$ freely and properly discontinuously by
  biholomorphisms preserving $t$.
\item[(b)] $U=\Ut/\Gamma$ is a complex threefold and $f=t\colon U\to D_\varepsilon$ is
  proper.
\item[(c)] $f^{-1}(0)\cong I_1\times E$, where $E=\C/(\Z\eb_2(0)+\Z\eb_3(0))$ and
  $I_1=t^{-1}(0)/\langle \nu\psi\rangle$ is a nodal rational curve. Near every point of
  $f^{-1}(0)$ there are local coordinates in which $f=xy$.
\item[(d)] $f^{-1}(D^*_\varepsilon)$ is the open subset
  $(\C^*_\zeta\times\C_w\times D^*_\varepsilon)/\Gamma$,
  on which $G_2,G_3,G_4$ act by
  $(\zeta,w)\mapsto(e^{\alpha_j}\zeta,w+\eb_j)$ and $(\zeta,w)\mapsto(\nu t\zeta,w+\eb_4)$.
  Its fibres are the tori with lattice
  \[\big\langle(2\pi,0),\ (-i\alpha_2,\eb_2),\ (-i\alpha_3,\eb_3),\ (-i\log(\nu t),\eb_4)\big\rangle.\]
\item[(e)] The section $t\mapsto[(\zeta{=}1,t),0]$ extends holomorphically over $t=0$.
\end{enumerate}
\end{proposition}

\begin{proof}
\emph{Step 0: constants.} Shrinking $\varepsilon_0$, we may assume the data are holomorphic
on a neighbourhood of $\overline D_{\varepsilon_0}$. Then there are constants
$c_0,C_0,C_1>0$ such that for $|t|\le\varepsilon_0$, real $a,b$ and $u\in\C$:
\begin{itemize}
\item $|a\eb_2(t)+b\eb_3(t)|\ge c_0(|a|+|b|)$;
\item $|\alpha_j(t)|+|\eb_4(t)|\le C_0|t|$;
\item $|\varphi_t(u)|\le C_1|t||u|$, where $\varphi_t\colon\C\to\R$ is the $\R$-linear
  functional with $\varphi_t(\eb_j(t))=\Rea\alpha_j(t)$ for $j=2,3$. It exists by (H2),
  and the estimate follows from (H1).
\end{itemize}
For $t\neq0$ put $P(t)=\log|\nu|+\log|t|-\varphi_t(\eb_4(t))$. Then $P(t)/\log|t|\to1$ as
$t\to0$. Choose $\varepsilon<\varepsilon_0$
such that $P(t)/\log|t|\in[\tfrac12,\tfrac32]$ for $0<|t|<\varepsilon$. This $\varepsilon$
depends only on the data, not on any compact set. For the two cusps used below, $P$ is
computed exactly in Remark~\ref{rem:Pexact}, and $\varepsilon=10^{-2}$ suffices.

\emph{Step 1: the action.}
\begin{itemize}
\item The maps are holomorphic and preserve $t$, since the $\C^*$-action and $t$ are
  holomorphic and $t\circ\psi=t$.
\item They commute, since $\psi$ commutes with the $\C^*$-action.
\item On $t^{-1}(0)\times\C$, $G_2^aG_3^bG_4^k$ acts by
  $(\nu^k\psi^k,\ w+a\eb_2(0)+b\eb_3(0))$, which is the identity only for $a=b=k=0$.
  Hence $\Gamma\cong\Z^3$.
\end{itemize}

\emph{Step 2: the open part.} For $t\ne0$ the fibre of $\cX$ is $\C^*_\zeta$ and $\psi$
acts by $\zeta\mapsto t\zeta$, which gives the formulas in (d). By
Lemma~\ref{lem:localdet} with $\ell=\log(\nu t)$, the lattice there has determinant
$-2\pi P(t)\,\omega(\eb_2,\eb_3)\neq0$.

\emph{Step 3: a normalised height.} On torus points of $\Ut$ put
$H(p,w)=\log|\zeta(p)|-\varphi_t(w)$. Since $\log|e^{\alpha_j}|=\Rea\alpha_j=\varphi_t(\eb_j)$,
\[
H\circ G_2=H,\qquad H\circ G_3=H,\qquad H\circ G_4=H+P(t).
\]

\emph{Step 4: freeness.} Suppose $\gamma=G_2^aG_3^bG_4^k$ fixes $(p,w)$.
\begin{itemize}
\item If $t(p)=0$, then $a\eb_2(0)+b\eb_3(0)=0$, so $a=b=0$, and $\nu^k\psi^k(p)=p$
  forces $k=0$, because $\nu^k\psi^k$ maps $D_n$ to $D_{n+k}$ and $P_n$ to $P_{n+k}$.
\item If $t(p)\ne0$, a fixed point means that $ae_2+be_3+ke_4\in\Z e_1$ for the
  lattice of Step 2, which is impossible unless $a=b=k=0$.
\end{itemize}

\emph{Step 5: proper discontinuity.} Let $K\subset\Ut$ be compact. Then
$K\subset\big(\bigcup_{|n|\le N}W_n(2)\big)\times\{|w|\le R\}$ for some $N,R$
(Lemma~\ref{lem:slabs}(1)). Let $\gamma=G_2^aG_3^bG_4^k$ and $q=(p,w)\in K$ with
$\gamma q\in K$.
\begin{itemize}
\item \emph{The $w$-coordinate.} Comparing $w$-coordinates gives
  $|a\eb_2+b\eb_3|\le 2R+|k|C_0\varepsilon$, so
  $|a|+|b|\le c_0^{-1}(2R+|k|C_0\varepsilon)$. It therefore suffices to bound $|k|$.
\item \emph{If $t(p)=0$.} $q$ and $\gamma q$ lie on components $D_m$ with $|m|\le N+1$,
  and $\gamma$ shifts component indices by $k$. So $|k|\le 2N+1$.
\item \emph{If $t(p)\neq0$.} By Lemma~\ref{lem:slabs}(5) and Step~0,
  \[\Big|\frac{H(q')}{\log|t|}\Big|\le N+2+\frac{C_1R\,\varepsilon}{\log(1/\varepsilon)}=:B_K\]
  for $q'\in\{q,\gamma q\}$. By Step~3, $H(\gamma q)-H(q)=kP(t)$, and
  $P(t)/\log|t|\ge\tfrac12$. Hence $|k|\le 4B_K$.
\end{itemize}

\emph{Step 6: properness of $f$.} Fix $\varepsilon'<\varepsilon$ and let
\[\Pi_t=\{s_2\eb_2(t)+s_3\eb_3(t):s_j\in[0,1]\},\qquad R_0=\max_{|t|\le\varepsilon'}\max_{u\in\Pi_t}|u|.\]
Let $q=(p,w)$ with $|t(p)|\le\varepsilon'$.
\begin{itemize}
\item If $t\neq0$: replace $q$ by $G_4^kq$ with $H\in(P(t),0]$ (Step 3), then by
  $G_2^aG_3^bq$ with $w\in\Pi_t$; $H$ is unchanged. Now
  $\log|\zeta|=H+\varphi_t(w)\in(\log|t|-C_2,\,C_2)$ with
  $C_2=|\log|\nu||+C_0C_1\varepsilon^2+C_1R_0\varepsilon$. So in the chart $U_0$,
  $|y_0|=|\zeta|<e^{C_2}$ and $|x_0|=|t|/|\zeta|<e^{C_2}$.
\item If $t=0$: use $G_4^k$ to move $p$ into $D_0$, then $G_2,G_3$ to put
  $w\in\Pi_0$.
\end{itemize}
Hence every $\Gamma$-orbit in $f^{-1}(\overline{D}_{\varepsilon'})$ meets the compact set
\[K_0=\Big\{(p,w):\ p\in\big(\{|x_0|,|y_0|\le e^{C_2}\}\cup\{|x_{-1}|,|y_{-1}|\le1\}\big),\
|t(p)|\le\varepsilon',\ w\in\Pi_{t(p)}\Big\}.\]
This uses $D_0\subset\{|x_0|=0,|y_0|\le1\}\cup\{y_{-1}=0,|x_{-1}|\le1\}$. So
$f^{-1}(\overline D_{\varepsilon'})$ is the image of $K_0$ and is compact.

\emph{Step 7: central fibre and local form.} At $t=0$, $G_2=(\mathrm{id},w+\eb_2(0))$,
$G_3=(\mathrm{id},w+\eb_3(0))$ and $G_4=(\nu\psi,w)$, by (H1). So
\[f^{-1}(0)=(t^{-1}(0)/\langle \nu\psi\rangle)\times E.\]
The quotient of the chain by $\nu\psi$ is $D_0$ with $P_{-1}\sim P_0$, i.e.\ $I_1$. In
the chart $U_n\times\C$, $f=x_ny_n$; this proves (c).

(b) follows from (a), Step 6 and the smoothness of $\Ut$.

(e) The point $(\zeta=1,t)$ has $x_0=t$ and $y_0=1$, which tends to the point
$(0,1)\in D_0\cap U_0$.
\end{proof}

\subsection{The cusp at $t=0$}
Take $\alpha_2=t$, $\alpha_3=-it$, $\eb_2=1-t^2$, $\eb_3=i(1+t^2)$, $\eb_4=2t$ and $\nu=1$.
Then (H1) holds, and (H2) holds with $\omega(\eb_2,\eb_3)=1-|t|^4$. Over $D^*_\varepsilon$ the
generators of Proposition~\ref{prop:filling}(d) are exactly $g_2,g_3,g_4$. Let
$U_0\to D_\varepsilon$ be the filling, and let
$\iota_0\colon U_0|_{D^*_\varepsilon}\xrightarrow{\ \sim\ }\Yo|_{D^*_\varepsilon}$ be the
identity of $(\C^*\times\C\times D^*_\varepsilon)/\Gamma^\circ$.

\subsection{The cusp at $t=\infty$}
\begin{lemma}\label{lem:atinfty}
On $\C^*\times\C\times\{|t|>1\}$ the map
\[\Theta(\zeta,w,t)=(\zeta e^{w/t},\ w/t^2,\ 1/t)=:(\zeta',w',s)\]
is a biholomorphism onto $\C^*\times\C\times D^*_1$, with inverse
$(\zeta',w',s)\mapsto(\zeta'e^{-w'/s},w'/s^2,1/s)$. It conjugates
\begin{align*}
g_2&\mapsto h_2=(e^{s}\zeta',\ w'+s^2-1),\qquad
g_3\mapsto h_3=(e^{is}\zeta',\ w'+i(1+s^2)),\\
g_4^{-1}&\mapsto h_4=(e^{-2}s\,\zeta',\ w'-2s).
\end{align*}
Hence $\Yo|_{|t|>1}\cong(\C^*\times\C\times D^*_1)/\langle h_2,h_3,h_4\rangle$. The data
\[
\alpha_2=s,\quad \alpha_3=is,\quad \eb_2=s^2-1,\quad \eb_3=i(1+s^2),\quad \eb_4=-2s,\quad \nu=e^{-2}
\]
satisfy (H1) and (H2), with $\omega(\eb_2,\eb_3)=|s|^4-1<0$.
\end{lemma}
\begin{proof}
For a map $(\zeta,w)\mapsto(m\zeta,w+u)$ with $m,u$ depending only on $t$,
\[\zeta e^{w/t}\mapsto me^{u/t}\cdot\zeta e^{w/t},\qquad w/t^2\mapsto w/t^2+u/t^2.\]
Insert $(m,u)=(e^t,1-t^2)$, $(e^{-it},i(1+t^2))$ and $(t,2t)$; see also
Appendix~\ref{app:checks}. The value of $\omega$ is a direct computation.
\end{proof}

Let $U_\infty\to D_\varepsilon$ ($s$-coordinate) be the filling for these data (after
shrinking $\varepsilon$ so that both fillings exist), and let
$\iota_\infty\colon U_\infty|_{D^*_\varepsilon}\xrightarrow{\sim}\Yo|_{|t|>1/\varepsilon}$ be
$\Theta^{-1}$. The central fibre is $I_1\times E_\infty$, where $E_\infty=\C/(\Z(-1)+\Z i)$ is
$\C/(\Z+i\Z)$ with the marked basis $(-1,i)$, of opposite orientation to the basis
$(1,i)$ at $0$.
The zero sections correspond, since $\Theta(1,0,t)=(1,0,1/t)$.

\begin{remark}[the heights $P$ exactly]\label{rem:Pexact}
By Lemma~\ref{lem:localdet}, $P=-\det_\R/(2\pi\,\omega(\eb_2,\eb_3))$, where $\det_\R$ is
the determinant of the lattice in Proposition~\ref{prop:filling}(d).
\begin{itemize}
\item At $t=0$ this lattice is $\Lambda_t$, and $\omega(\eb_2,\eb_3)=1-r^4$ with $r=|t|$.
  So Lemma~\ref{lem:det-family} gives
  \[P_0(r)=\log r-\frac{2r^2}{1-r^2}.\]
\item At $\infty$, put $\rho=|s|$. In the additive coordinates $(z'',w')$ with
  $\zeta'=e^{iz''}$, the map $\Theta$ is the $\C$-linear map
  $(z,w)\mapsto(z-iw/t,\ w/t^2)$. Its real determinant is $|t|^{-4}=\rho^4$. It maps
  $(e_1,e_2,e_3,-e_4)$ to a basis of the lattice at $\infty$, up to adding multiples of
  $e_1$ to the basis vectors.
\item Hence $\det_\R=-\rho^4\cdot2\pi(1+r^2)(2r^2-(1-r^2)\log r)$ with $r=1/\rho$. This
  equals $-2\pi(1+\rho^2)(2-(1-\rho^2)\log\rho)$. With $\omega(\eb_2,\eb_3)=\rho^4-1$
  this gives
  \[P_\infty(\rho)=\log\rho-\frac{2}{1-\rho^2}.\]
\end{itemize}
Therefore, for $0<r<1$,
\[\frac{P_0(r)}{\log r}=1+\frac{2r^2}{(1-r^2)\,|\log r|},\qquad
\frac{P_\infty(r)}{\log r}=1+\frac{2}{(1-r^2)\,|\log r|}.\]
Both are increasing in $r$ and tend to $1$ as $r\to0$. So for $0<r<10^{-2}$ they lie in
$(1,1.000044)$ and $(1,1.434339)$, respectively, and $\varepsilon=10^{-2}$ satisfies
the condition of Step~0 at both cusps. At $r=1/16$ the second ratio is $1.724\ldots>\tfrac32$,
so $\varepsilon=1/16$ does not work at $\infty$.
\end{remark}

%======================================================================
\section{The logarithmic transform and the threefold $Y$}\label{sec:twist}
%======================================================================

The following map is a logarithmic transform of multiplicity one in the sense of
Kodaira \cite{Kodaira63}: a regluing by a holomorphic section over a punctured disc.
Its topological type is a nonzero class $[e_4]$ in the coinvariants
$\Lambda/(T-1)\Lambda$.

\begin{lemma}\label{lem:Phi}
On $\Yo|_{D^*_\varepsilon}$ define
\[
\Phi[\zeta,w,t]=\big[\,e^{i\pi s(s-1)}\zeta,\ \ w+2ts,\ \ t\,\big],\qquad s=\tfrac{\log t}{2\pi i}.
\]
\begin{enumerate}
\item $\Phi$ is a well-defined biholomorphism over the identity of $D^*_\varepsilon$ and
  commutes with $g_2,g_3,g_4$.
\item In additive coordinates $\Phi$ is fibrewise translation by
  $\tilde\sigma(s)=(\pi s(s-1),\,2ts)$, and $\tilde\sigma(s+1)-\tilde\sigma(s)=e_4(s)$,
  where $e_4(s)=(2\pi s,2t)$.
\end{enumerate}
\end{lemma}
\begin{proof}
Replacing $s$ by $s+1$ multiplies $e^{i\pi s(s-1)}$ by
\[e^{i\pi((s+1)s-s(s-1))}=e^{2\pi is}=t\]
and adds $2t$ to $w+2ts$. So the two values differ by $g_4$, and $\Phi$ is well
defined on the quotient. It is locally given by a holomorphic formula. It commutes
with the $g_j$ because all of them are of the form $(\zeta,w)\mapsto(m\zeta,w+u)$ with
$m,u$ depending only on $t$. Its inverse is given by the same formula with signs
reversed. Statement (2) is a direct computation, using $-i\log t=2\pi s$.
\end{proof}

\begin{construction}\label{con:Y}
Fix $\varepsilon$ as above, and put $Y_{\rm mid}=\Yo|_{\varepsilon/2<|t|<2/\varepsilon}$. Define
\[
Y=\big(U_0\ \sqcup\ Y_{\rm mid}\ \sqcup\ U_\infty\big)\big/\!\sim
\]
with the identifications
\begin{itemize}
\item $q\sim\Phi(\iota_0(q))$ for $q\in U_0$ with $\varepsilon/2<|t(q)|<\varepsilon$;
\item $q\sim\iota_\infty(q)$ for $q\in U_\infty$ with $\varepsilon/2<|s(q)|<\varepsilon$.
\end{itemize}
The maps $t$ and $s=1/t$ glue to $f\colon Y\to\PP^1$.
\end{construction}

\begin{proposition}\label{prop:Y}
$Y$ is a compact connected complex manifold of dimension $3$, and $f$ is a proper
surjective holomorphic map. Over $\C^*$ its fibres are the tori $X_t$, and
$f^{-1}(0)\cong I_1\times E_0$, $f^{-1}(\infty)\cong I_1\times E_\infty$.
\end{proposition}
\begin{proof}
The gluing maps are biholomorphisms between open subsets, compatible with $f$ and
equal to the identity on the base.
\begin{itemize}
\item \emph{Hausdorff.} Each of $U_0$, $Y_{\rm mid}$ and $U_\infty$ is the full
  preimage of an open subset of $\PP^1$, and maps injectively onto an open subset of
  $Y$. The gluing maps cover the identity of the base, so every fibre of $f$ lies
  entirely in at least one piece. Let $q_1\ne q_2$.
  \begin{itemize}
  \item If $f(q_1)\ne f(q_2)$, separate them by preimages of disjoint discs.
  \item If $f(q_1)=f(q_2)=t_0$, choose a piece containing the whole fibre over $t_0$.
    It contains both points, and they can be separated inside it.
  \end{itemize}
\item \emph{Compact.} $Y$ is the union of the compact sets
  $f^{-1}(|t|\le3\varepsilon/4)$, $f^{-1}(3\varepsilon/4\le|t|\le4/(3\varepsilon))$ and
  $f^{-1}(|s|\le3\varepsilon/4)$, by Proposition~\ref{prop:filling}(b) and
  Lemma~\ref{lem:Yo}.
\item \emph{Connected.} Each piece is connected and the pieces overlap.
\end{itemize}
\end{proof}

%======================================================================
\section{The torus action}\label{sec:torus}
%======================================================================

\begin{proposition}\label{prop:torus}
For $c=(c_2,c_3)\in\R^2$ define, on each piece:
\begin{align*}
\text{on }\Yo:\quad &A_c[\zeta,w,t]=[e^{(c_2-ic_3)t}\zeta,\ w+c_2(1-t^2)+ic_3(1+t^2),\ t],\\
\text{on }\Ut_0:\quad &\tilde A_c(p,w)=(e^{(c_2-ic_3)t}\cdot p,\ w+c_2\eb_2(t)+c_3\eb_3(t)),\\
\text{on }\Ut_\infty:\quad &\tilde A'_c(p,w')=(e^{(c_2+ic_3)s}\cdot p,\ w'+c_2(s^2-1)+ic_3(1+s^2)).
\end{align*}
These define an action of $T^2=\R^2/\Z^2$ on $Y$ by biholomorphisms, preserving $f$,
which is free and real-analytic. Allowing $c\in\C^2$ gives a holomorphic action of
$(\C^*)^2=\C^2/\Z^2$.
\end{proposition}
We write $m_c$ for the multiplier of the $\cX$-factor: $m_c=e^{(c_2-ic_3)t}$ on $\Ut_0$
and $m_c=e^{(c_2+ic_3)s}$ on $\Ut_\infty$. On both pieces, then,
$\tilde A_c(p,w)=(m_c\cdot p,\ w+c_2\eb_2+c_3\eb_3)$, with the data $\eb_2,\eb_3$ of the
respective cusp.
\begin{proof}
\emph{Well defined on each piece.}
\begin{itemize}
\item Each formula is additive in $c$, and $c=(1,0),(0,1)$ give $g_2,g_3$ (resp.\
  $G_2,G_3$ and $h_2,h_3$). So $\Z^2$ acts trivially on the quotients.
\item Each formula is a map $(\zeta,w)\mapsto(m\zeta,w+u)$ with $m,u$ depending only on
  the base, so it commutes with $g_4$, $G_4$, $h_4$ and with $\Phi$.
\end{itemize}
\emph{Compatible with the gluings.} On $U_0|_{D^*}$ the formula is that of $\Yo$, since
$\iota_0$ is the identity. The conjugate of $A_c$ by $\Theta$ is $\tilde A'_c$, by the
computation in the proof of Lemma~\ref{lem:atinfty} with
$(m,u)=(e^{(c_2-ic_3)t},c_2(1-t^2)+ic_3(1+t^2))$:
\[me^{u/t}=e^{(c_2+ic_3)/t},\qquad u/t^2=c_2(s^2-1)+ic_3(1+s^2).\]
(See Appendix~\ref{app:checks}.)

\emph{Free.}
\begin{itemize}
\item On $X_t$, $t\in\C^*$, $A_c$ is translation by $c_2e_2(t)+c_3e_3(t)$. This lies in
  $\Lambda_t$ only if $c\in\Z^2$, by $\R$-independence of $e_1,\dots,e_4$.
\item On $f^{-1}(0)=I_1\times E_0$ the action is $(\mathrm{id},\ u\mapsto u+c_2+ic_3)$,
  since $e^{(c_2-ic_3)\cdot0}=1$, and this is free on $T^2$.
\item On $f^{-1}(\infty)$ it is translation by $-c_2+ic_3$ on $E_\infty$.
\end{itemize}
\emph{Holomorphic.} Each $A_c$ is holomorphic in all variables for fixed $c$, and also
for complex $c$.
\end{proof}

\begin{corollary}\label{cor:W}
$W:=Y/T^2$ is a smooth, closed, connected, orientable $4$-manifold, and $\pi\colon Y\to W$
is a smooth principal $T^2$-bundle. The map $f$ factors as $f=f_W\circ\pi$ with
$f_W\colon W\to\PP^1$ smooth and proper.
\end{corollary}
\begin{proof}
A free smooth action of a compact Lie group gives a smooth principal bundle. The
quotient of an oriented manifold by a free action of a connected group is orientable.
The function $f$ is $T^2$-invariant.
\end{proof}

%======================================================================
\section{Topology of $W$}\label{sec:topology}
%======================================================================

Let $N_0=\pi(U_0)$, $M=\pi(Y_{\rm mid})$ and $N_\infty=\pi(U_\infty)$. These are open in
$W$ and cover it. We have $N_0\cap N_\infty=\emptyset$, and $O_0=N_0\cap M$, $O_\infty=N_\infty\cap M$ lie over
the annuli $\varepsilon/2<|t|<\varepsilon$ and $1/\varepsilon<|t|<2/\varepsilon$.

\subsection{The regular part}
Let $e_1^*,\dots,e_4^*$ be the real dual basis of $(e_1,\dots,e_4)$ on $\C^2=\R^4$. Under
$e_4\mapsto e_4+e_1$ the dual basis changes to $(e_1^*-e_4^*,e_2^*,e_3^*,e_4^*)$, so
$e_4^*$ is a well-defined continuous function of $t\in\C^*$. Moreover
$e_4^*(\Lambda_t)\subset\Z$ and $e_4^*(e_2)=e_4^*(e_3)=0$.

\begin{lemma}\label{lem:regular}
\begin{enumerate}
\item For $t\in\C^*$, $f_W^{-1}(t)=X_t/T^2=\C^2/(\Lambda_t+\R e_2+\R e_3)$ is a $2$-torus
  with $H_1=\Z a\oplus\Z b$, where $a=[e_1]$ and $b=[e_4]$.
\item $M\to\{\varepsilon/2<|t|<2/\varepsilon\}$ is a smooth $T^2$-bundle with monodromy
  $a\mapsto a$, $b\mapsto b+a$. With base point on the image of the zero section and
  $x$ the loop $\theta\mapsto[1,0,r_0e^{i\theta}]$,
  \[\pi_1(M)=\langle a,b,x\mid [a,b],\ [x,a],\ xbx^{-1}=b a^{\pm1}\rangle.\]
\item The map $\Psi\colon M\to S^1\times S^1$, $\Psi[v,t]=(\arg t,\ e_4^*(v)\bmod\Z)$, is well
  defined and continuous. It induces an isomorphism
  $\pi_1(M)/\langle\!\langle a\rangle\!\rangle\cong\Z^2$ with $x\mapsto(1,0)$, $b\mapsto(0,1)$.
\end{enumerate}
\end{lemma}
\begin{proof}
(1) $e_1,\dots,e_4$ is a real basis, so $\R^4/(\R e_2+\R e_3)\cong\R e_1\oplus\R e_4$.

(2) Ehresmann's theorem applies, since $f_W$ is a proper submersion over $\C^*$. The
monodromy is that of Lemma~\ref{lem:Yo}, and the zero section gives the semidirect
product description.

(3) $\Psi$ is well defined by the properties of $e_4^*$. Modulo the normal closure of
$a$, $\pi_1(M)$ is the free abelian group on $b,x$, and $\Psi_*$ sends these to the
standard basis.
\end{proof}

\subsection{The local pieces}
\begin{lemma}\label{lem:Nlocal}
Let $j\in\{0,\infty\}$.
\begin{enumerate}
\item $N_j\cong\cX^\varepsilon/\langle\hat\psi\rangle$, where
  $\hat\psi(p)=\kappa(t(p))\cdot \nu\,\psi(p)$, $\kappa=\exp\big((\gamma_2-i\gamma_3)t\big)$
  (resp.\ $\exp((\gamma_2+i\gamma_3)s)$), and $(\gamma_2,\gamma_3)\colon D_\varepsilon\to\R^2$ is the
  real-analytic solution of $\gamma_2\eb_2+\gamma_3\eb_3=-\eb_4$. The action of
  $\langle\hat\psi\rangle$ is free and properly discontinuous.
\item $\pi_1(N_j)\cong\Z$, generated by $\hat\psi$, and
  $H_q(N_j;\Z)\cong\Z,\Z,\Z,0,0$ for $q=0,1,2,3,4$. In particular $\chi(N_j)=1$.
\item Let $a_j,b_j,x_j$ generate $\pi_1(O_j)$: the fibre classes $a_j=[e_1]$ and
  $b_j=[e_4]$ (resp.\ $[-e_4]$), and the loop $x_j$ around the cusp on the zero section
  of $U_j$. Under $O_j\subset N_j$ these map to $1$, $\hat\psi$ and $1$.
\end{enumerate}
\end{lemma}
\begin{proof}
(1) $\R^2$ acts on $\Ut_j$ by $\tilde A_c$, and $\Gamma\subset\R^2\times\langle G_4\rangle$, so
$N_j=U_j/T^2=\Ut_j/(\R^2\times\langle G_4\rangle)$.
\begin{itemize}
\item Since $\eb_2(t),\eb_3(t)$ is a real basis of $\C$, the slice $\{w=0\}$ meets every
  $\R^2$-orbit exactly once. The map
  $(p,w)\mapsto m_c\cdot p$, with $c=c(t,w)$ solving $w+c_2\eb_2+c_3\eb_3=0$, is a
  real-analytic submersion identifying $\Ut_j/\R^2$ with $\cX^\varepsilon$. (Here
  $m_c=e^{(c_2-ic_3)t}$ for $j=0$ and $m_c=e^{(c_2+ic_3)s}$ for $j=\infty$.)
\item $G_4$ maps $(p,0)$ to $(\nu\psi(p),\eb_4)$, which lies on the $\R^2$-orbit of
  $(\hat\psi(p),0)$.
\end{itemize}
Since $\R^2$ is central in $\R^2\times\langle G_4\rangle$,
$N_j=(\Ut_j/\R^2)/\langle G_4\rangle=\cX^\varepsilon/\langle\hat\psi\rangle$. Explicitly, the
slice inclusion $p\mapsto(p,0)$ and the slice map $(p,w)\mapsto m_{c(t,w)}\cdot p$ are
smooth and equivariant. They induce mutually inverse smooth maps
$\cX^\varepsilon/\langle\hat\psi\rangle\leftrightarrow U_j/T^2$, since $U_j\to U_j/T^2$ is a
smooth principal bundle. So these spaces are diffeomorphic, and under the
diffeomorphism $f_W$ corresponds to $t$.
\begin{itemize}
\item \emph{Free.} Suppose $\hat\psi^k$ fixes $p$. Then $\tilde A_cG_4^k$ fixes $(p,0)$ for
  some $c$. Its image in $U_j$ is fixed by the class $[c]\in T^2$, so $c\in\Z^2$ by
  freeness of the $T^2$-action. Then $\tilde A_cG_4^k\in\Gamma$, and $k=0$ by freeness of
  $\Gamma$.
\item \emph{Properly discontinuous.} Let $K\subset\cX^\varepsilon$ be compact. Put
  $\widetilde K=\tilde A_{[0,1]^2}(K\times\{0\})$, which is compact. If
  $\hat\psi^kK\cap K\ne\emptyset$, write $\tilde A_cG_4^k=G_2^{\lfloor c_2\rfloor}G_3^{\lfloor c_3\rfloor}G_4^k\circ
  \tilde A_{\{c\}}$ with $\{c\}\in[0,1)^2$; the element
  $G_2^{\lfloor c_2\rfloor}G_3^{\lfloor c_3\rfloor}G_4^k\in\Gamma$ then moves $\widetilde K$ to meet
  $\widetilde K$. By Proposition~\ref{prop:filling}(a) only finitely many $k$ occur.
\end{itemize}

(2) By Lemma~\ref{lem:Xtop}, $\cX^\varepsilon$ is simply connected, so it is the universal
cover of $N_j$ and $\pi_1(N_j)=\langle\hat\psi\rangle\cong\Z$. Now $\hat\psi=\kappa'\cdot\psi$
with $\kappa'=\nu\,\kappa(t)$, which has the continuous logarithm $\log\nu+(\gamma_2\mp i\gamma_3)t$ (sign $-$ at $0$, $+$ at $\infty$, with $s$ in place of $t$). So Lemma~\ref{lem:Xtop} gives $\hat\psi_*=\psi_*=\tau$ on homology. The
Wang sequence (F2) gives short exact sequences
\[0\to\mathrm{coker}(\tau_*-1\mid H_q(\cX^\varepsilon))\to H_q(N_j)\to\ker(\tau_*-1\mid H_{q-1}(\cX^\varepsilon))\to0.\]
Since $\tau_*=1$ on $H_0=\Z$, $H_1=0$, and $\tau-1$ is injective with cokernel $\Z$ on
$\Z[\tau^{\pm1}]$, these give
\begin{itemize}
\item $H_0(N_j)=\Z$;
\item $H_1(N_j)=\ker(\tau_*-1\mid H_0)=\Z$;
\item $H_2(N_j)=\mathrm{coker}(\tau-1\mid\Z[\tau^{\pm1}])=\Z$;
\item $H_3(N_j)=\ker(\tau-1\mid\Z[\tau^{\pm1}])=0$, and $H_4(N_j)=0$.
\end{itemize}

(3) The loops $a_j$ (the $\zeta$-circle) and $x_j$ (the loop $t=re^{i\theta}$,
$\zeta=1$, $w=0$; on it $c=0$ in the slice map) lift to closed loops in
$\cX^\varepsilon$. The fibre loop $b_j$ lifts to a path from $p$ to $\hat\psi(p)$.
\end{proof}

\subsection{The effect of the logarithmic transform}
\begin{lemma}\label{lem:twistclass}
Let $\gamma$ be the image in $M$, via $\Phi\circ\iota_0$, of the loop $x_0$. Then
$\Psi_*(\gamma)=(1,1)$, hence $\gamma\equiv xb\pmod{\langle\!\langle a\rangle\!\rangle}$. On the other
hand, $x_\infty$ maps to $x^{-1}$ under $\iota_\infty$.
\end{lemma}
\begin{proof}
The loop is $\theta\mapsto[\tilde\sigma(s(\theta)),t(\theta)]$ in additive coordinates, with
$t(\theta)=re^{i\theta}$ and $s(\theta)$ a continuous branch. Its $\arg t$ winds once.
Since $e_4^*$ is continuous and single-valued, the change of $e_4^*$ along the loop is
\[e_4^*\big(\tilde\sigma(s(0)+1)-\tilde\sigma(s(0))\big)=e_4^*(e_4)=1\]
by Lemma~\ref{lem:Phi}(2). The loop $x$ has $e_4^*\equiv0$ along it. For $\infty$:
$\iota_\infty$ maps zero section to zero section (\S\ref{sec:filling}), and a
counterclockwise loop in $s$ is clockwise in $t$.
\end{proof}

\begin{proposition}\label{prop:pi1W}
$\pi_1(W)=1$.
\end{proposition}
\begin{proof}
All pieces and intersections are path-connected, so we apply van Kampen to
$W=(N_0\cup M)\cup N_\infty$.
\begin{itemize}
\item By Lemma~\ref{lem:Nlocal}(3), gluing $N_\infty$ imposes $a=1$ and $x_\infty=1$,
  i.e.\ $x=1$ by Lemma~\ref{lem:twistclass}. It
  also identifies the generator of $\pi_1(N_\infty)$ with $b^{\mp1}$, up to a factor
  in $\langle\!\langle a\rangle\!\rangle$.
\item Gluing $N_0$ along $O_0$ via $\Phi\circ\iota_0$ imposes $a=1$ and $\gamma=1$. It
  also identifies the generator of $\pi_1(N_0)$ with $b$. (Basepoint changes only
  conjugate these relations.)
\end{itemize}
By Lemma~\ref{lem:twistclass}, $\gamma=xb\,n$ with $n\in\langle\!\langle a\rangle\!\rangle$. So
$1=\gamma=b$ in the quotient. Thus all generators $a,b,x$ of $\pi_1(M)$ and the generators
of $\pi_1(N_j)$ are trivial.
\end{proof}

\begin{proposition}\label{prop:chiW}
$\chi(W)=2$.
\end{proposition}
\begin{proof}
All pieces have finitely generated homology, so the Mayer--Vietoris sequence of
the open cover $\{N_0,\,M,\,N_\infty\}$ gives
\[\chi(W)=\chi(N_0)+\chi(M)+\chi(N_\infty)-\chi(O_0)-\chi(O_\infty).\]
$M$, $O_0$ and $O_\infty$ are $T^2$-bundles over annuli, so their Euler characteristics
vanish; $\chi(N_j)=1$ by Lemma~\ref{lem:Nlocal}(2).
\end{proof}

\begin{corollary}\label{cor:homotopysphere}
$W$ is a smooth homotopy $4$-sphere.
\end{corollary}
\begin{proof}
$W$ is closed, oriented and simply connected.
\begin{itemize}
\item Hence $H_1=0$, $H_3\cong H^1=0$, $H_4=\Z$, and $H_2\cong H^2\cong\mathrm{Hom}(H_2,\Z)$
  is free.
\item Then $\chi(W)=2+b_2=2$, so $H_2=0$.
\item The collapse map $W\to S^4$ of degree one is a homology isomorphism of simply
  connected spaces, hence a homotopy equivalence by Whitehead's theorem.
\end{itemize}
\end{proof}

%======================================================================
\section{Proof of the Main Theorem}\label{sec:proof}
%======================================================================

\begin{proof}[Proof of the Main Theorem]
\emph{The bundle is trivial.} By Corollary~\ref{cor:W}, $Y\to W$ is a principal
$T^2$-bundle. By Corollary~\ref{cor:homotopysphere}, $H^2(W;\Z^2)=0$. By (F1) the bundle
is smoothly trivial, so $Y\cong W\times T^2$.

\emph{$h$-cobordism.} By (F3), $W$ is $h$-cobordant to $S^4$: both are simply connected
with zero intersection form. Let $H^5$ be such an $h$-cobordism. Then $H\times T^2$ is a
$7$-dimensional $h$-cobordism between $W\times T^2$ and $S^4\times T^2$, with fundamental
group $\Z^2$.

\emph{$s$-cobordism.} Its Whitehead torsion lies in $\mathrm{Wh}(\Z^2)=0$ (F4). By the
$s$-cobordism theorem (F4),
\[Y\cong W\times T^2\cong S^4\times T^2.\]

\emph{Transport.} If $F\colon S^4\times T^2\to Y$ is a diffeomorphism and $J_Y$ is the
complex structure of $Y$, then $F^*J_Y$ is an almost complex structure on
$S^4\times T^2$. It is integrable, since $F$ is a biholomorphism from
$(S^4\times T^2,F^*J_Y)$ to $Y$; equivalently, $N_{F^*J_Y}=F^*N_{J_Y}=0$.
\end{proof}

\subsection{The quotient $W$ is the standard $S^4$}
The Main Theorem only needs $W$ to be a homotopy sphere. Using Iwase's classification
of torus fibrations with twin singular fibres \cite{Iwase84}, one can say more. We
follow the definitions of \cite[\S2]{Matsumoto85}. A \emph{(good) torus fibration} is a
proper surjective smooth map $f\colon M^4\to B^2$ of oriented manifolds with two
properties. First, near every point it is given in local complex coordinates by
$\xi=z_1^mz_2^n$ or $\xi=\bar z_1^mz_2^n$. Second, it is a smooth $T^2$-bundle off a
discrete set of critical values. A singular fibre of type $I_1^\pm$ is an immersed
$2$-sphere with one transverse self-intersection of sign $\pm1$. A \emph{twin} singular
fibre ($Tw$) consists of two embedded $2$-spheres meeting transversely in two points,
of opposite signs.

\begin{proposition}\label{prop:Wstandard}
$W$ is diffeomorphic to $S^4$. Consequently $Y\cong S^4\times T^2$ as principal
$T^2$-bundles, and the diffeomorphism of the Main Theorem can be chosen
$T^2$-equivariant.
\end{proposition}
\begin{proof}
Fix an orientation of $W$.

\emph{Step 1: $f_W$ is a torus fibration with two singular fibres, each of type
$I_1^+$ or $I_1^-$.}
\begin{itemize}
\item Over $\C^*$, $f_W$ is a smooth $T^2$-bundle (Lemma~\ref{lem:regular}).
\item By the proof of Lemma~\ref{lem:Nlocal}(1), there is a diffeomorphism
  $N_j\cong\cX^\varepsilon/\langle\hat\psi\rangle$ under which $f_W$ corresponds to $t$
  (to $s$ at $\infty$). The quotient map
  $\cX^\varepsilon\to\cX^\varepsilon/\langle\hat\psi\rangle$ is a covering. So the holomorphic
  chart $(x_0,y_0)$ of $U_0\subset\cX^\varepsilon$ near $P_0$ gives smooth complex
  coordinates on $W$ near the singular point of $f_W^{-1}(0)$, in which $f_W=x_0y_0$.
\item If these coordinates are positively oriented for $W$, then $f_W=z_1z_2$ with
  $(z_1,z_2)=(x_0,y_0)$. Otherwise $f_W=\bar z_1z_2$ with the positively oriented
  coordinates $(z_1,z_2)=(\bar x_0,y_0)$.
\item Elsewhere on the singular fibre, $dt\ne0$, so $f_W$ is locally $\xi=z_1$.
\item The singular fibre is the image of the chain $t^{-1}(0)$ modulo $\hat\psi$,
  i.e.\ $D_0$ with $P_{-1}$ and $P_0$ identified. This is an immersed $2$-sphere with
  one transverse self-intersection. It is not multiple, since $f_W=x_0y_0$ is reduced.
  The same holds at $\infty$.
\end{itemize}

\emph{Step 2: merging the nodes.} Let $D\subset\PP^1$ be a closed disc containing $0$
and $\infty$ in its interior and no other critical value, e.g.\ the complement of a
small open disc around $t=1$. The monodromy around $\partial D$ is trivial, since
$\partial D$ also bounds a disc without critical values. By
\cite[Thm.~4.1]{Matsumoto85}, the fibration $f_W|_{f_W^{-1}(D)}$ can be deformed,
without changing it near the boundary, into a torus fibration with a single singular
fibre, of type $Tw$ with divisor $R+S$. This gives a torus fibration $W\to S^2$ whose
only singular fibre is a non-multiple twin fibre.

\emph{Step 3: Iwase's theorem.} Iwase proved the following (see
\cite[\S1]{Matsumoto85} for this statement).
\begin{quote}
Let $M\to S^2$ be a torus fibration whose singular fibres are non-multiple twin
fibres, and assume $e(M)\neq0$. Then the diffeomorphism type of $M$ is determined by
$\pi_1(M)$, $e(M)$, $w_2(M)$ and the parity of the intersection form on $H_2(M;\Z)$.
\end{quote}
The $4$-sphere is the total space of such a fibration, with a single twin fibre
\cite{Matsumoto82}, \cite[\S2, \S4]{Matsumoto85}. Both $W$ and $S^4$ have $\pi_1=1$ and
$e=2$ (Propositions~\ref{prop:pi1W} and \ref{prop:chiW}). Both have $H^2(\,\cdot\,;\Z)=0$
and $H^2(\,\cdot\,;\Z/2)=0$ (Corollary~\ref{cor:homotopysphere}), so $w_2=0$ and the
intersection form is zero, hence even. Therefore $W\cong S^4$.

The last statement follows from $Y\cong W\times T^2$ (\S\ref{sec:proof}).
\end{proof}

\begin{remark}
Proposition~\ref{prop:Wstandard} identifies the total space $W$ with $S^4$. It does not
show that $f_W$ is equivalent, as a fibration, to Matsumoto's fibration of $S^4$
\cite{Matsumoto82,Matsumoto85,GompfStipsicz99}, although the two have the same singular
fibres, monodromy and twist class. The proposition gives a second proof of the Main
Theorem, which replaces (F3)--(F4) by Iwase's classification. The proof in
\S\ref{sec:proof} uses only the classical results (F1)--(F4).
\end{remark}

%======================================================================
\section{Further properties (Theorem B)}\label{sec:props}
%======================================================================

\begin{proof}[Proof of Theorem B]
(1) This is Proposition~\ref{prop:Y}.

(2) The $T^2$-statements are Proposition~\ref{prop:torus} and \S\ref{sec:proof}. For
$(z_0,w_0)\in\C^2$, the formulas
\[
[\zeta,w,t]\mapsto[e^{iz_0}\zeta,w+w_0,t],\qquad (p,w)\mapsto(e^{iz_0}\cdot p,w+w_0),\qquad
(p,w')\mapsto(e^{iz_0+w_0s}\cdot p,\ w'+w_0s^2)
\]
on $\Yo$, $\Ut_0$ and $\Ut_\infty$ commute with all generators and with $\Phi$. They are
compatible with $\Theta$, by the same computation as in Lemma~\ref{lem:atinfty}. So
they define a holomorphic $\C^2$-action, which is transitive on each $X_t$.

\emph{The kernel.} If $(z_0,w_0)$ acts trivially, then $(z_0,w_0)\in\Lambda_t$ for every $t$.
Writing $(z_0,w_0)=\sum n_je_j(t)$ with locally constant integers, the second
coordinate $n_2(1-t^2)+in_3(1+t^2)+2n_4t$ must be constant in $t$, so
$n_2=n_3=n_4=0$. Hence the kernel is $\Z(2\pi,0)$, and
$\Aut^0(Y)\supset\C^2/\Z(2\pi,0)\cong\C^*\times\C$.

(3) \emph{$c_1=0$.} Let $O\subset X_t$ be a $T^2$-orbit. Under $Y\cong W\times T^2$ it is
$\{pt\}\times T^2$, which generates $H_2(Y;\Z)\cong\Z$. There is an exact sequence of
holomorphic vector bundles on $X_t$,
\[0\to TX_t\to TY|_{X_t}\to N_{X_t}\to0,\qquad N_{X_t}=f^*T_t\PP^1|_{X_t}\cong\cO_{X_t}.\]
Here $TX_t\cong\cO^2$ is trivialised by translations. By additivity of $c_1$,
$c_1(TY|_{X_t})=0$. So $c_1(Y)\cdot[O]=0$, and $c_1(Y)=0$ because
$H^2(Y;\Z)\cong\Z$ is free.

The sequence splits as a sequence of smooth complex vector bundles, but not
holomorphically. Its extension class is the Kodaira--Spencer class of $f$ at $t$. That
class is nonzero, because the period ratio $p_{13}/p_{12}=i(1+t^2)/(1-t^2)$ of
Lemma~\ref{lem:periods} has derivative $4it/(1-t^2)^2\ne0$; at $t=\pm1$ use
$p_{12}/p_{13}$ instead. Only $c_1$ is used above.

\emph{The canonical bundle.} Consider the meromorphic $3$-form $\Omega$ given on $\Yo$ by
$dz\wedge dw\wedge dt=\tfrac1{i}\tfrac{d\zeta}{\zeta}\wedge dw\wedge dt$.
\begin{itemize}
\item It is invariant under the $g_j$ and $\Phi$, which are of the form
  $(z,w,t)\mapsto(z+\mu(t),w+u(t),t)$.
\item On $\Ut_0$, since $\tfrac{d\zeta}{\zeta}\wedge dw\wedge dt=-\tfrac{d\zeta}{\zeta}\wedge dt\wedge dw$,
  it equals $\tfrac{1}{i}\,dx_n\wedge dy_n\wedge dw$ (Lemma~\ref{lem:toric}(6)), which is
  holomorphic and nowhere zero.
\item Near $\infty$, $\Omega=-s^{-4}\,dz''\wedge dw'\wedge ds$ with $\zeta'=e^{iz''}$
  (Appendix~\ref{app:checks}), and $dz''\wedge dw'\wedge ds$ is a nowhere-zero
  holomorphic form on $\Ut_\infty$.
\end{itemize}
Hence $\mathrm{div}(\Omega)=-4F_\infty$, and $K_Y\cong\cO_Y(-4F_\infty)=f^*\cO(-4)$, because
$F_\infty=f^*[\infty]$ is reduced. If $f^*\cO(-4m)$ were trivial for some $m\ge1$, it
would have a nonzero section. But $f_*\cO_Y=\cO_{\PP^1}$ (connected fibres) gives
$H^0(Y,f^*\cO(-4m))=H^0(\PP^1,\cO(-4m))=0$.

(4) Under $Y\cong W\times T^2$, a smooth fibre is $\pi^{-1}(f_W^{-1}(t))=f_W^{-1}(t)\times T^2$,
and $H_2(W)=0$, so $[F]=0$ in $H_4(Y;\Z)$.
\begin{itemize}
\item \emph{Not balanced.} A balanced metric $\omega_Y$ ($d\omega_Y^2=0$) would give
  $0<\int_F\omega_Y^2=\langle[\omega_Y^2],[F]\rangle=0$.
\item \emph{Not K\"ahler.} K\"ahler metrics are balanced. Alternatively,
  $H^2(Y;\R)=\R a$ with $a^2=0$.
\end{itemize}
(As a consistency check: $c_1(Y)=-c_1(f^*\cO(-4))=4\,\mathrm{PD}[F]=0$.)

(5) Pulling back rational functions by $f$ gives $f^*\C(\PP^1)\subset\mathcal M(Y)$.
Conversely, let $h\in\mathcal M(Y)$, and let $\mathcal E$ be the countable set of
Proposition~\ref{prop:NS}.
\begin{itemize}
\item \emph{Over a disc.} Let $B\subset\C^*$ be an open disc. Pull $h$ back along the
  holomorphic covering $\C^2\times B\to f^{-1}(B)$ to a meromorphic function $H$,
  invariant under $\Lambda_t$.
\item \emph{Restriction to fibres.} The irreducible components of the polar set of $H$
  form a locally finite family. Every fibre $\C^2\times\{t\}$ meets the compact set
  $\{0\}\times K$ for $t$ in a compact $K\subset B$. So only a discrete set of fibres
  lies in the polar set.
\item For every other $t$, $H(\cdot,t)$ is a $\Lambda_t$-periodic meromorphic function
  on $\C^2$, i.e.\ a meromorphic function on $X_t$. For $t\notin\mathcal E$ it is
  constant by Proposition~\ref{prop:NS}.
\item \emph{$H$ depends only on $t$.} The set of such $t$ is dense in $B$. So
  $\partial_zH$ and $\partial_wH$ vanish on a dense subset of the open set where they
  are holomorphic, hence identically. Since the complement of the polar set is
  connected, $H=c_B\circ\mathrm{pr}_B$ for a meromorphic function $c_B$ on $B$. (Take
  $c_B=H(z_0,w_0,\cdot)$ for $(z_0,w_0)$ such that $\{(z_0,w_0)\}\times B$ is not
  contained in the polar set.)
\item \emph{Gluing.} The functions $c_B$ glue to a meromorphic function $c$ on $\C^*$
  with $h=f^*c$ on $\Yo$.
\item \emph{The ends.} Let $\sigma_0$ be the holomorphic section of $U_0\to
  D_\varepsilon$ of Proposition~\ref{prop:filling}(e). Off a discrete set,
  $\sigma_0(t)$ is a point where $h=f^*c$ is holomorphic. So $\sigma_0(D_\varepsilon)$
  is not contained in the polar set of $h$, and $h\circ\sigma_0$ is meromorphic on
  $D_\varepsilon$. It equals $c$ on $D_\varepsilon^*$, so $c$ is meromorphic at $0$. The
  same argument with the section of $U_\infty$ applies at $\infty$.
\end{itemize}
Thus $c\in\C(\PP^1)$, and $h=f^*c$, since both sides are meromorphic on $Y$ and agree on
the dense open set $\Yo$. Hence $\mathcal M(Y)=f^*\C(\PP^1)$, and $a(Y)=1$.
\end{proof}

\begin{remark}[deformations]\label{rem:deform}
Replace $g_2,g_3,g_4$ by
\[(e^{\lambda t}\zeta,w+1-t^2),\qquad(e^{-i\lambda t}\zeta,w+i(1+t^2)),\qquad
(e^{i\beta}t\zeta,w+2t)\]
with $\lambda>0$ and $\beta\in\C$. Then the determinant of Lemma~\ref{lem:det-family}
becomes $2\pi(1+r^2)D(r)$, where $r=|t|$, $b=\Ima\beta$ and
\[D(r)=2\lambda r^2+(b-\ln r)(1-r^2).\]

\emph{The admissible region.} $D$ is positive for every $r>0$ if and only if
\begin{equation}\label{eq:admissible}
|\Ima\beta-\lambda|<\sqrt{\lambda(\lambda+1)}+\operatorname{arsinh}\sqrt\lambda.
\end{equation}
\begin{itemize}
\item $D(1)=2\lambda>0$.
\item For $0<r<1$, $D(r)>0$ if and only if $b>g(r)$, where
  $g(r)=\ln r-2\lambda r^2/(1-r^2)$. The function $g$ tends to $-\infty$ at both ends
  of $(0,1)$. Its derivative $1/r-4\lambda r/(1-r^2)^2$ vanishes only where
  $1-r^2=2\sqrt\lambda\,r$, i.e.\ at $r_\lambda=\sqrt{\lambda+1}-\sqrt\lambda$. There
  $r_\lambda^2/(1-r_\lambda^2)=r_\lambda/(2\sqrt\lambda)$, so
  \[S_\lambda:=\max_{(0,1)}g=\ln r_\lambda-\sqrt\lambda\,r_\lambda
  =\lambda-\sqrt{\lambda(\lambda+1)}-\operatorname{arsinh}\sqrt\lambda.\]
\item For $r>1$, substitute $r=1/\rho$. Since $1/(1-\rho^2)=1+\rho^2/(1-\rho^2)$,
  $D(1/\rho)>0$ if and only if $b<2\lambda-g(\rho)$.
\item So $D>0$ on $(0,\infty)$ if and only if $S_\lambda<b<2\lambda-S_\lambda$, which is
  \eqref{eq:admissible}.
\end{itemize}
The region \eqref{eq:admissible} contains $\Ima\beta=0$ for every $\lambda>0$, since
$\sqrt{\lambda(\lambda+1)}>\lambda$. It is not all of $\{\lambda>0\}\times\C$. For
example, at $\lambda=1$ and $\beta=-2i$ one has $D(\tfrac12)=-1+\tfrac34\ln2<0$. So
$D$ vanishes for some $r\in(\tfrac12,1)$, and the four period vectors fail to span a
lattice there. Outside \eqref{eq:admissible} the construction breaks down.

\emph{Modified proofs.} For $(\lambda,\beta)$ satisfying \eqref{eq:admissible}, the
proofs above go through with the following changes:
\begin{itemize}
\item use $\Theta_\lambda(\zeta,w,t)=(\zeta e^{\lambda w/t},w/t^2,1/t)$ at $\infty$;
\item there the filling data are $\alpha_2=\lambda s$, $\alpha_3=i\lambda s$, and the
  constant becomes $\nu=e^{-i\beta-2\lambda}$;
\item the twist becomes translation by $(\pi s(s-1)+\beta s,\ 2ts)$;
\item the $T^2$-action has multipliers $e^{\lambda(c_2-ic_3)t}$ near $0$ and
  $e^{\lambda(c_2+ic_3)s}$ near $\infty$.
\end{itemize}
This gives a family of complex structures on $S^4\times T^2$, parametrised by the
region \eqref{eq:admissible}. The original construction is $(\lambda,\beta)=(1,0)$. We
have not determined which members are biholomorphic.
\end{remark}

%======================================================================
\section{Remarks and consistency checks}\label{sec:remarks}
%======================================================================

\begin{remark}[Leray spectral sequence: a heuristic check]
This remark is not used in the proof; the Betti numbers follow rigorously from
$Y\cong W\times T^2$.
\begin{itemize}
\item Rationally, the monodromy-invariant covectors are $e_2^*,e_3^*,e_4^*$.
\item In the Leray spectral sequence of $f$, the twist class $[e_4]$ gives
  $d_2(e_4^*)=\pm\omega_{\PP^1}$, while $d_2(e_2^*)=d_2(e_3^*)=0$.
\item This yields $b_1(Y)=2$, and the class of $H^2(\PP^1)$ dies.
\item The surviving class $e_2^*\wedge e_3^*$ accounts for $b_2(Y)=1$.
\end{itemize}
A complete Leray computation would need the sheaves $R^qf_*\Q$ at the singular fibres;
we do not carry it out.
\end{remark}

\begin{remark}[signature]
Matsumoto's local signature of a genus-one node is $\mp\tfrac23$ for positive and
negative nodes. Since $\sigma(W)=0$, the two nodes of $f_W$ must have opposite signs,
in accordance with Remark~\ref{rem:chirality}.
\end{remark}

\begin{remark}[polarised families have no cusps over $\C^*$]\label{rem:polarised}
Let $\{A_t\}_{t\in\C^*}$ be a holomorphic family of polarised abelian surfaces.
\begin{itemize}
\item Its period map lifts to the universal cover, $\C\to\mathfrak H_2$. Siegel space
  is biholomorphic to a bounded domain, so the lift is constant by Liouville's theorem.
\item Hence the monodromy fixes a point of $\mathfrak H_2$ and has finite order. After a
  finite base change the family is trivial near $0$ and $\infty$; in particular it has
  no cusp there.
\end{itemize}
So the family of \S\ref{sec:family}, whose monodromy $T$ has infinite order, admits
no polarisation over $\C^*$. Its period map takes values in the space of complex
structures on $\R^4$, which is not hyperbolic, and is non-constant.

This argument concerns polarisations of the whole family. It does not by itself
decide whether individual fibres are projective, and some are: $X_1$ is an abelian
surface. Non-algebraicity of the very general fibre is a statement about integral
Hodge classes. We prove it directly from the periods in Proposition~\ref{prop:NS}:
$\mathrm{NS}(X_t)=0$ outside a countable set. The same non-algebraic phenomenon
underlies \cite{AC26}.
\end{remark}

\begin{remark}[a necessary condition]\label{rem:necessity}
Let $f\colon Y\to\PP^1$ be a fibration by complex $2$-tori with finitely many singular
fibres and connected fibres, with $Y\simeq S^4\times T^2$. Assume that the restriction
$H^0(\PP^1,R^1f_*\Q)\to H^1(F_t;\Q)^{\pi_1}$ to a smooth fibre is injective. A section
that vanishes on one smooth fibre vanishes on the whole regular locus, so this holds
whenever each singular fibre has a neighbourhood $N$ with $H^1(N;\Q)\to H^1(F_t;\Q)$
injective. That is the case for the fibres $I_1\times E$ used here, since
$H_1(F_t)\to H_1(I_1\times E)$ is onto. The assumption is used in the second and
third steps below.
\begin{itemize}
\item The base contributes no $H^1$, so $b_1(Y)=2$ means
  $\ker(d_2\colon H^0(R^1)\to H^2(\PP^1;\Q))$ has rank $2$.
\item If $d_2=0$, then $f^*[\mathrm{pt}]\ne0$ in $H^2$, and the cup product of the two
  classes in $H^1$ is a second, independent class (it projects to
  $\psi_1\wedge\psi_2\neq0$ in $E_\infty^{0,2}$). That gives $b_2\ge2$, a contradiction.
\item The other ranks give $b_1\ne2$, so $\mathrm{rank}\,H^0(R^1)=3$. Hence every
  monodromy fixes a common rational hyperplane $\Pi\subset H^1(F_t;\Q)$ (or acts
  trivially).
\end{itemize}
Monodromies of a holomorphic family preserve the complex orientation of the fibres,
so they have determinant $1$. An element $T$ of $SL_4$ fixing $\Pi$ pointwise
satisfies $T^*=I+u\otimes\ell$ with $\ell|_\Pi=0$ and $\det T=1+\ell(u)=1$. So $u\in\Pi$,
and $T$ is a transvection. Dually, every monodromy is a transvection
$\lambda\mapsto\lambda+u_\gamma(\lambda)\lambda_0$ with a common $\lambda_0$. A transvection
$I+N$ with $N^2=0$ satisfies $(I+N)^k=I+kN$. So nontrivial monodromies of finite
order, as in the $(3,4,\infty)$ family, cannot occur. In our construction
$\lambda_0=e_1$.
\end{remark}

\begin{remark}[questions]
\begin{enumerate}
\item $W$ is diffeomorphic to $S^4$ (Proposition~\ref{prop:Wstandard}). Is $f_W$
  equivalent to Matsumoto's fibration? Can one write an explicit diffeomorphism
  $W\cong S^4$, and hence an explicit integrable $J$ on the round $S^4\times T^2$?
\item Compute the Dolbeault cohomology and the Kuranishi space of $Y$.
\item Twist by $n[e_4]$ instead of $[e_4]$, and call the quotient $W_n$. The arguments
  of \S\ref{sec:topology} give $\chi(W_n)=2$, and $\pi_1(W_n)=\Z/n$ for $n\ne0$. So for
  $n\ne0$, $W_n$ is a rational homology $4$-sphere. Untwisted, $\pi_1(W_0)=\Z$, so
  $b_1=b_3=1$ and $b_2=2$. Steps 1--2 of the proof of
  Proposition~\ref{prop:Wstandard} apply to every $W_n$. So Iwase's theorem
  determines $W_n$ from $\pi_1(W_n)$, $w_2(W_n)$ and the parity of its intersection
  form. Which $4$-manifolds are they?
\item Calabi's question for $V^2$ of genus $\ge2$ remains open as far as we know. For
  the products of spheres of dimensions $10$ and $12$ that remain open, see
  \S\ref{sec:higherdim}.
\end{enumerate}
\end{remark}

%======================================================================
\section{Products of spheres}\label{sec:spheres}
%======================================================================

In this section we combine the Main Theorem with the 2026 constructions and determine
which products of spheres of dimension at most eight admit complex structures. The one
new ingredient besides the Main Theorem is Theorem~C. It applies a Calabi--Eckmann type
construction to the complex $S^2\times S^4$ of Liang, Liu and Zhang \cite{LLZ26}, and
gives complex structures on $S^4\times S^3\times S^{2q+1}$.

Unlike the Main Theorem, Theorem~C and Corollary~D depend on recent preprints. We list
these inputs explicitly in \S\ref{sec:inputs}.

Throughout this section, products of spheres carry their standard smooth structures, and
$S^n$ always has $n\ge1$. Products of spheres of dimension $n$ correspond to partitions of
$n$.

\subsection{Inputs}\label{sec:inputs}
Besides classical constructions (complex projective spaces, elliptic curves, Hopf and
Calabi--Eckmann manifolds, and products), we use:
\begin{enumerate}
\item[(I1)] $S^6$ admits a complex structure \cite{AC26,Engel26}. We write
  $X_{\rm sph}\cong S^6$ for the compact complex threefold of \cite{AC26}.
\item[(I2)] $S^2\times S^4$ admits a complex structure \cite{GFL26,LLZ26}.
\item[(I3)] \emph{The package (LLZ).} There is a compact connected complex threefold
  $X_+$, a diffeomorphism $X_+\cong S^2\times S^4$, and a holomorphic line bundle
  $L_+\to X_+$ whose first Chern class generates $H^2(X_+;\Z)\cong\Z$ \cite{LLZ26}.
\end{enumerate}
These inputs depend on one another as follows.
\begin{itemize}
\item (I3) implies (I2).
\item (I3) assumes more than (I1). It starts from the specific threefold $X_{\rm sph}$
  together with the following properties:
  \begin{itemize}
  \item $\mathrm{Aut}^0(X_{\rm sph})\cong\C^*$;
  \item the fixed curve of this action has normal weights $\pm1$;
  \item the structure of the cusp fibre, with its $\C^*$-invariant double curves.
  \end{itemize}
  See \cite[Thm.~4.5, Prop.~4.6, Props.~9.23--9.24, Rem.~9.25]{AC26}.
\item Engel's independent account \cite{Engel26} verifies the construction of
  $X_{\rm sph}$, but treats these auxiliary properties only in a remark.
\item (I2) also rests on \cite{AC26} by either route: \cite{GFL26} reuse its analytic
  framework, and \cite{LLZ26} start from $X_{\rm sph}$.
\end{itemize}
So none of (I1)--(I3) is independent of \cite{AC26}.

\subsection{Elliptic bundles over a complex base}\label{sec:CE}
The following proposition is the Calabi--Eckmann construction \cite{CalabiEckmann53},
with $\PP^p\times\PP^q$ replaced by an arbitrary compact complex base. It is standard;
see \cite{Morimoto63,Hofer93,GoldsteinProkushkin04,GGP08} for this and related
constructions on torus bundles. We give a complete proof, since everything below rests
on it.

\begin{proposition}\label{prop:CE}
Let $B$ be a compact connected complex manifold of dimension $n$, let $L_1,L_2\to B$ be
holomorphic line bundles, and fix $\tau\in\C\setminus\R$. Let
\[
\mathcal P=L_1^\times\times_BL_2^\times=\{(v_1,v_2)\in L_1\oplus L_2:\ v_1\neq0\neq v_2\},
\]
and let $\C$ act on $\mathcal P$ by
\[
t\cdot(v_1,v_2)=(e^{t}v_1,\ e^{\tau t}v_2).
\]
Then:
\begin{enumerate}
\item $Z=\mathcal P/\C$ is a compact complex manifold of dimension $n+1$;
\item $Z\to B$ is a holomorphic fibre bundle with fibre the elliptic curve
  $E_\tau=\C/(\Z+\tau\Z)$, and structure group $E_\tau$ acting by translations;
\item for any hermitian metrics on $L_1$ and $L_2$, $Z$ is diffeomorphic to the
  principal $T^2$-bundle
  \[\mathcal Q=S(L_1)\times_BS(L_2)=\{(v_1,v_2)\in\mathcal P:\ |v_1|=|v_2|=1\}.\]
\end{enumerate}
\end{proposition}

The proof occupies Lemmas~\ref{lem:Htau}--\ref{lem:Q}. Let $G=(\C^*)^2$. The map
$\exp\colon\C^2\to G$, $(z_1,z_2)\mapsto(e^{z_1},e^{z_2})$, is a surjective holomorphic
homomorphism with kernel $2\pi i\Z^2$. Put
\[H_\tau=\{(e^t,e^{\tau t}):t\in\C\}=\exp\big(\C\cdot(1,\tau)\big)\subset G.\]

\begin{lemma}\label{lem:Htau}
\begin{enumerate}
\item The map $t\mapsto(e^t,e^{\tau t})$ is injective. So the $\C$-action on $\mathcal P$
  is the action of the subgroup $H_\tau\cong\C$, acting through the fibrewise
  $G$-action $(w_1,w_2)\cdot(v_1,v_2)=(w_1v_1,w_2v_2)$.
\item The map
  \[\rho\colon G\to E'_\tau=\C/2\pi i(\Z+\tau\Z),\qquad (w_1,w_2)\mapsto[\log w_2-\tau\log w_1],\]
  is a well-defined surjective holomorphic homomorphism with kernel $H_\tau$. Hence $H_\tau$
  is a closed complex subgroup, and $G/H_\tau\cong E'_\tau$ is an elliptic curve,
  isomorphic to $E_\tau$ via $z\mapsto z/2\pi i$.
\end{enumerate}
\end{lemma}
\begin{proof}
(1) Suppose $e^t=e^{\tau t}=1$. Then $t=2\pi im$ and $\tau t=2\pi in$ with
$m,n\in\Z$, so $\tau m=n$. Since $\tau\notin\R$, this forces $m=0$, hence $t=0$.

(2) \emph{Well defined.} Changing the branches of the logarithms by $2\pi im$ and $2\pi in$
changes $\log w_2-\tau\log w_1$ by $2\pi i(n-\tau m)$. So $\rho$ is well defined. It is
locally given by a holomorphic formula, it is a homomorphism, and it is surjective since
$\rho(1,e^z)=[z]$.

\emph{Kernel.} Clearly $\rho(e^t,e^{\tau t})=[\tau t-\tau t]=0$. Conversely, suppose
$\rho(w)=0$. Choose logarithms $(z_1,z_2)$ of $w$ with $z_2-\tau z_1=2\pi i(n-\tau m)$. Then
$(z_1-2\pi im,\ z_2-2\pi in)$ is also a logarithm of $w$, and it satisfies
$z_2'=\tau z_1'$. So $w=(e^{t},e^{\tau t})$ with $t=z_1-2\pi im$.

\emph{Quotient.} Since $1$ and $\tau$ are $\R$-independent, $2\pi i(\Z+\tau\Z)$ is a lattice
in $\C$, so $E'_\tau$ is a compact complex torus of dimension one. A holomorphic
homomorphism of complex Lie groups has constant rank. If it is surjective, the rank is
maximal by Sard's theorem, so it is a submersion. Hence $\ker\rho=H_\tau$ is a closed
complex subgroup, and $\rho$ induces an isomorphism $G/H_\tau\cong E'_\tau$ of complex Lie
groups.
\end{proof}

\begin{lemma}\label{lem:Z}
$Z=\mathcal P/H_\tau$ is a compact complex manifold of dimension $n+1$, and $Z\to B$ is a
holomorphic fibre bundle with fibre $E'_\tau$ and structure group $E'_\tau$ acting by
translations. The projection $\mathcal P\to Z$ is a holomorphic surjective submersion
whose fibres are exactly the $H_\tau$-orbits.
\end{lemma}
\begin{proof}
\emph{Trivialisations of $\mathcal P$.} Choose an open cover $\{U_\alpha\}$ of $B$ and
nowhere-zero holomorphic sections $s^{(j)}_\alpha$ of $L_j|_{U_\alpha}$. Write
$s^{(j)}_\beta=g^{(j)}_{\alpha\beta}\,s^{(j)}_\alpha$ with $g^{(j)}_{\alpha\beta}\in\cO^*(U_{\alpha\beta})$, and put
$g_{\alpha\beta}=(g^{(1)}_{\alpha\beta},g^{(2)}_{\alpha\beta})\colon U_{\alpha\beta}\to G$. The maps
\[
U_\alpha\times G\to\mathcal P|_{U_\alpha},\qquad (b,w)\mapsto\big(w_1s^{(1)}_\alpha(b),\ w_2s^{(2)}_\alpha(b)\big),
\]
are biholomorphisms. Their transition maps are $(b,w)\mapsto(b,g_{\alpha\beta}(b)\,w)$, so
$\mathcal P$ is a holomorphic principal $G$-bundle. In these charts $H_\tau$ acts by
multiplication on the $G$-factor. This commutes with the transition maps because $G$ is
abelian.

\emph{The glued manifold.} Glue the complex manifolds $U_\alpha\times E'_\tau$ along
$U_{\alpha\beta}\times E'_\tau$ by the biholomorphisms
\[(b,e)\mapsto\big(b,\ e+\rho(g_{\alpha\beta}(b))\big).\]
They satisfy the cocycle condition, because $g_{\alpha\beta}$ does and $\rho$ is a homomorphism.
Call the result $Z'$. It is a complex manifold and a holomorphic fibre bundle over $B$ with
fibre $E'_\tau$ and translation transition maps.
\begin{itemize}
\item \emph{Hausdorff.} Points in different fibres are separated by preimages of
  disjoint open sets of $B$. Two points in the fibre over $b$ lie in a single chart
  $U_\alpha\times E'_\tau$ with $b\in U_\alpha$, and are separated there.
\item \emph{Second countable.} The cover may be taken countable.
\item \emph{Compact.} Cover $B$ by finitely many compact sets $K_\alpha\subset U_\alpha$. Then
  $Z'$ is the union of the compact images of the sets $K_\alpha\times E'_\tau$.
\end{itemize}

\emph{Identification with $\mathcal P/H_\tau$.} The local maps
$U_\alpha\times G\to U_\alpha\times E'_\tau$, $(b,w)\mapsto(b,\rho(w))$, are compatible with
the transition maps. So they glue to a holomorphic map $\pi\colon\mathcal P\to Z'$.
\begin{itemize}
\item $\pi$ is a surjective submersion, because $\rho$ is.
\item Its fibres are the $H_\tau$-orbits, because $\ker\rho=H_\tau$.
\end{itemize}
So $\pi$ induces a continuous bijection $\mathcal P/H_\tau\to Z'$. It is open, since
submersions are open maps, hence a homeomorphism. We identify $Z=\mathcal P/H_\tau$ with
$Z'$, and use Lemma~\ref{lem:Htau} to identify $E'_\tau$ with $E_\tau$.
\end{proof}

Fix hermitian metrics on $L_1$ and $L_2$, and define
\[\mu\colon\mathcal P\to\R^2,\qquad \mu(v)=(\log|v_1|,\ \log|v_2|),\]
so that $\mathcal Q=\mu^{-1}(0)$. Let $\mathrm A_\tau\colon\C\to\R^2$ be the $\R$-linear
map $\mathrm A_\tau(t)=(\Rea t,\ \Rea(\tau t))$. For $t=a+ib$,
$\mathrm A_\tau(t)=(a,\ a\Rea\tau-b\Ima\tau)$, so $\det\mathrm A_\tau=-\Ima\tau\neq0$ and
$\mathrm A_\tau$ is an isomorphism. Since $|e^tv_1|=e^{\Rea t}|v_1|$ and
$|e^{\tau t}v_2|=e^{\Rea(\tau t)}|v_2|$,
\begin{equation}\label{eq:mu}
\mu(t\cdot v)=\mu(v)+\mathrm A_\tau(t)\qquad(t\in\C,\ v\in\mathcal P).
\end{equation}

\begin{lemma}\label{lem:Q}
$\mathcal Q$ is a compact smooth submanifold of $\mathcal P$ and a principal
$T^2$-bundle over $B$. The composite $\Phi\colon\mathcal Q\hookrightarrow\mathcal P\to Z$
is a diffeomorphism.
\end{lemma}
\begin{proof}
\emph{$\mathcal Q$.} By \eqref{eq:mu}, the derivative of $\mu$ along each $H_\tau$-orbit is
the isomorphism $\mathrm A_\tau$. So $\mu$ is a submersion, and $\mathcal Q=\mu^{-1}(0)$ is
a smooth submanifold. It is the fibre product of the unit circle bundles $S(L_1)$ and
$S(L_2)$. Hence it is a principal $U(1)^2$-bundle over $B$, and compact because $B$ is.

\emph{Injectivity of $\Phi$.} If $v,v'\in\mathcal Q$ and $v'=t\cdot v$, then by
\eqref{eq:mu} $\mathrm A_\tau(t)=\mu(v')-\mu(v)=0$, so $t=0$ and $v'=v$.

\emph{The inverse.} Define $\tilde\Psi\colon\mathcal P\to\mathcal Q$ by
\[\tilde\Psi(v)=t(v)\cdot v,\qquad t(v)=-\mathrm A_\tau^{-1}\mu(v).\]
By \eqref{eq:mu}, $\mu(\tilde\Psi(v))=0$, so $\tilde\Psi$ does land in $\mathcal Q$, and
it is smooth. It is $H_\tau$-invariant: \eqref{eq:mu} gives $t(s\cdot v)=t(v)-s$, so
$\tilde\Psi(s\cdot v)=t(v)\cdot v=\tilde\Psi(v)$. Since $\mathcal P\to Z$ is a surjective
submersion (Lemma~\ref{lem:Z}), $\tilde\Psi$ descends to a smooth map $\Psi\colon Z\to\mathcal Q$.
\begin{itemize}
\item For $v\in\mathcal Q$ we have $t(v)=0$, so $\Psi\circ\Phi=\mathrm{id}_{\mathcal Q}$.
\item $\tilde\Psi(v)$ lies in the $H_\tau$-orbit of $v$, so $\Phi\circ\Psi=\mathrm{id}_Z$.
\end{itemize}
\end{proof}

Lemmas~\ref{lem:Htau}, \ref{lem:Z} and~\ref{lem:Q} prove Proposition~\ref{prop:CE}.

\begin{remark}[one line bundle]\label{rem:onebundle}
If $L_2$ is trivial, then $Z\cong L_1^\times/\lambda^{\Z}$, the quotient of $L_1^\times$ by
fibrewise multiplication by $\lambda=e^{2\pi i/\tau}$.
\begin{itemize}
\item The slice $\{v_2=1\}$ meets every $H_\tau$-orbit.
\item The elements of $H_\tau$ preserving the slice are those with $e^{\tau t}=1$, i.e.\
  $t\in\frac{2\pi i}{\tau}\Z$. They act on $v_1$ by powers of $\lambda$.
\item $|\lambda|=\exp(2\pi\Ima\tau/|\tau|^2)\neq1$.
\end{itemize}
Conversely, every $\lambda$ with $|\lambda|\neq1$ arises, from $\tau=2\pi i/\log\lambda$. In
this case Proposition~\ref{prop:CE}(3) says $L_1^\times/\lambda^{\Z}\cong S(L_1)\times S^1$.
For $L_1=\cO_{\PP^{p}}(-1)$ this is the Hopf manifold
$(\C^{p+1}\setminus0)/\lambda^{\Z}\cong S^{2p+1}\times S^1$.
\end{remark}

\subsection{Circle bundles}\label{sec:circlebundles}
We use the following standard facts.
\begin{enumerate}
\item[(C1)] Every complex line bundle $L$ over a manifold has a hermitian metric. The
  unit circle bundle $S(L)$ is a smooth principal $U(1)$-bundle with associated line
  bundle $L$. Its Chern class is $c_1(L)$, the Euler class of the underlying oriented
  real plane bundle. Rescaling the metric gives an isomorphic circle bundle.
\item[(C2)] Smooth principal $U(1)$-bundles over a manifold $B$ are classified up to
  smooth isomorphism by their Chern class in $H^2(B;\Z)$ \cite{Husemoller}.
\item[(C3)] If $\varphi\colon B'\to B$ is a diffeomorphism and $\pi\colon\mathcal Q\to B$ is a
  smooth fibre bundle, then $(b',x)\mapsto x$ is a diffeomorphism from the total space of
  $\varphi^*\mathcal Q$ onto $\mathcal Q$.
\item[(C4)] Give $\cO(-1)\to\PP^q$ the metric induced from $\C^{q+1}$. Then
  $S(\cO(-1))=\{(\ell,v):v\in\ell,\ |v|=1\}$, and $(\ell,v)\mapsto v$ is a diffeomorphism
  onto $S^{2q+1}$. For $q=1$ this is the Hopf fibration.
  \begin{itemize}
  \item $S(\cO(1))$ is $S(\cO(-1))$ with the inverse $U(1)$-action, so over $\PP^1$
    both total spaces are $S^3$.
  \item For $q=0$, $\PP^0$ is a point and $S(\cO(-1))=S^1$.
  \end{itemize}
\end{enumerate}

\begin{lemma}\label{lem:productcircle}
Let $B=B_1\times B_2$, and let $L_j=\mathrm{pr}_j^*L'_j$ for line bundles $L'_j\to B_j$,
with the pulled-back metrics. Then $S(L_1)\times_BS(L_2)\cong S(L'_1)\times S(L'_2)$.
\end{lemma}
\begin{proof}
$S(L_1)=S(L'_1)\times B_2$ and $S(L_2)=B_1\times S(L'_2)$. A pair
$((u_1,b_2),(b_1,u_2))$ lies in the fibre product exactly when $u_1$ lies over $b_1$ and
$u_2$ lies over $b_2$. So $((u_1,b_2),(b_1,u_2))\mapsto(u_1,u_2)$ is a diffeomorphism onto
$S(L'_1)\times S(L'_2)$, with inverse $(u_1,u_2)\mapsto((u_1,\pi(u_2)),(\pi(u_1),u_2))$.
\end{proof}

\begin{lemma}\label{lem:S3S4}
Let $X$ be a manifold diffeomorphic to $S^2\times S^4$, and let $L\to X$ be a complex line
bundle whose Chern class generates $H^2(X;\Z)$. Then $S(L)$ is diffeomorphic to
$S^3\times S^4$.
\end{lemma}
\begin{proof}
Let $\varphi\colon S^2\times S^4\to X$ be a diffeomorphism.
\begin{itemize}
\item By Künneth, $H^2(S^2\times S^4;\Z)=\Z x$ with $x=\mathrm{pr}_1^*a$, where
  $a=c_1(\cO_{\PP^1}(1))$ under $S^2=\PP^1$. Since $\varphi^*$ is an isomorphism,
  $\varphi^*c_1(L)=\varepsilon x$ for some $\varepsilon=\pm1$.
\item The circle bundle $\mathrm{pr}_1^*S(\cO(\varepsilon))$ over $S^2\times S^4$ has Chern
  class $\varepsilon x$. By (C2), $\varphi^*S(L)\cong\mathrm{pr}_1^*S(\cO(\varepsilon))$ as smooth
  principal bundles.
\item The total space of $\mathrm{pr}_1^*S(\cO(\varepsilon))$ is
  $S(\cO(\varepsilon))\times S^4\cong S^3\times S^4$, by (C4).
\item By (C3), the total space of $\varphi^*S(L)$ is diffeomorphic to $S(L)$.
\end{itemize}
\end{proof}

\subsection{The complex $S^2\times S^4$ of Liang--Liu--Zhang}\label{sec:LLZ}
We recall the construction of \cite{LLZ26}, with references for each step. We then give
an independent proof that $c_1(L_+)$ is primitive, which is the property Theorem~C uses.
\begin{enumerate}
\item \emph{The source.} The threefold $X_{\rm sph}$ of \cite{AC26} has
  $\mathrm{Aut}^0(X_{\rm sph})\cong\C^*$. The $\C^*$-action fixes a smooth rational curve $C$
  with normal weights $\pm1$. One of the double curves $D_-$ of the cusp fibre is
  $\C^*$-invariant and meets $C$ in two points $Q,Q'$ \cite[\S1]{LLZ26}.
\item \emph{Blow-up and quotient.} Let $\widetilde X=\mathrm{Bl}_C(X_{\rm sph})$, with
  exceptional divisor $\mathcal E$.
  \begin{itemize}
  \item The involution $-1\in\C^*$ lifts to an involution $\tilde\iota$ of $\widetilde X$.
    In local coordinates $\tilde\iota(z,x,t)=(z,-x,t)$, with $\mathcal E=\{x=0\}$, so
    $\tilde\iota$ fixes $\mathcal E$ pointwise and acts freely elsewhere.
  \item The quotient has coordinates $(z,x^2,t)$. So
    $\varpi\colon\widetilde X\to X_-=\widetilde X/\langle\tilde\iota\rangle$ is a holomorphic
    double cover of a compact complex manifold, branched along the smooth divisor
    $\mathcal F=\varpi(\mathcal E)$, and $\varpi|_{\mathcal E}\colon\mathcal E\to\mathcal F$ is a
    biholomorphism \cite[\S2.1]{LLZ26}.
  \end{itemize}
\item \emph{The line bundle on $X_-$.} The anti-invariant part of $\varpi_*\cO_{\widetilde X}$
  is a line bundle $L_-^{-1}$. One has $L_-^{\otimes2}\cong\cO_{X_-}(\mathcal F)$
  \cite[\S4.1]{LLZ26} and
  \begin{equation}\label{eq:qL}
  \varpi^*L_-\cong\cO_{\widetilde X}(\mathcal E)\qquad\text{\cite[(4)]{LLZ26}}.
  \end{equation}
\item \emph{The flop.} The curve $C_-=\varpi(\widetilde D_-)$, the image of the strict
  transform of $D_-$, is a smooth rational curve with normal bundle
  $\cO(-1)\oplus\cO(-1)$ \cite[Lemma 3.1]{LLZ26}. The analytic Atiyah flop along $C_-$
  gives a compact complex threefold $X_+$ with flopped curve $C_+$
  \cite[Def.~3.2]{LLZ26}. Put $X^\circ=X_-\setminus C_-=X_+\setminus C_+$.
\item \emph{The line bundle on $X_+$.} Let $\widehat X$ be the common blow-up, with
  exceptional divisor $R\cong\PP^1\times\PP^1$. There is a unique holomorphic line bundle
  $L_+$ on $X_+$ with
  $\mathrm{Bl}_{C_+}^*L_+\cong\mathrm{Bl}_{C_-}^*L_-\otimes\cO_{\widehat X}(R)$
  \cite[(6)]{LLZ26}. Since $R$ lies over $C_\pm$, restricting to the complement of $R$ gives
  \begin{equation}\label{eq:LU}
  L_+|_{X^\circ}\cong L_-|_{X^\circ}.
  \end{equation}
\item \emph{Topology.} $X_+$ is simply connected, with torsion-free homology and
  $H^2(X_+;\Z)\cong\Z$ \cite[Prop.~3.5]{LLZ26}. By Wall's and Jupp's classification
  \cite{Wall66,Jupp73} there is an orientation-preserving diffeomorphism
  $X_+\cong S^2\times S^4$ \cite[Thm.~1.2 and \S5.2]{LLZ26}.
\end{enumerate}

\cite{LLZ26} show that $c_1(L_+)$ generates $H^2(X_+)$ by computing
$c_1(L_-)^3=1$ \cite[Lemma 4.1]{LLZ26} and using
$H^2(X_-)\cong H^2(X^\circ)\cong H^2(X_+)$. The following direct argument avoids the cubic
form.

\begin{lemma}\label{lem:primitive}
$c_1(L_+)$ generates $H^2(X_+;\Z)$. More precisely, there is a smooth rational curve
$\Sigma\subset X^\circ$ with $\deg(L_+|_\Sigma)=-1$.
\end{lemma}
\begin{proof}
\emph{Choice of $\Sigma$.} The curve $D_-$ is not contained in $C$, so $D_-\cap C$ is finite;
it is $\{Q,Q'\}$. Choose $c\in C\setminus D_-$. Let $\Sigma_c=\mathrm{Bl}_C^{-1}(c)\cong\PP^1$ be
the fibre of $\mathcal E\to C$ over $c$, and put $\Sigma=\varpi(\Sigma_c)$.

\emph{$\Sigma$ lies in $X^\circ$.}
\begin{itemize}
\item The strict transform $\widetilde D_-$ meets $\mathcal E$ only over $D_-\cap C$, so
  $\Sigma_c\cap\widetilde D_-=\emptyset$.
\item $\tilde\iota$ fixes $\Sigma_c\subset\mathcal E$ pointwise, so
  $\Sigma_c\cap\tilde\iota(\widetilde D_-)=\tilde\iota(\Sigma_c\cap\widetilde D_-)=\emptyset$.
\item Since $\varpi^{-1}(C_-)=\widetilde D_-\cup\tilde\iota(\widetilde D_-)$, it follows that
  $\Sigma\cap C_-=\emptyset$, i.e.\ $\Sigma\subset X^\circ$.
\end{itemize}

\emph{The degree.} $\varpi|_{\mathcal E}$ is a biholomorphism onto $\mathcal F$, so
$\varpi|_{\Sigma_c}\colon\Sigma_c\to\Sigma$ is an isomorphism of smooth rational curves.
Using \eqref{eq:LU} and \eqref{eq:qL},
\[
\deg(L_+|_\Sigma)=\deg(L_-|_\Sigma)=\deg\big((\varpi^*L_-)|_{\Sigma_c}\big)=\deg\big(\cO_{\widetilde X}(\mathcal E)|_{\Sigma_c}\big)=-1.
\]
The last equality holds because $\cO_{\widetilde X}(\mathcal E)|_{\mathcal E}$ is the
tautological bundle $\cO_{\PP(N_{C})}(-1)$, which has degree $-1$ on each fibre of
$\mathcal E\to C$.

\emph{Conclusion.} So $\langle c_1(L_+),[\Sigma]\rangle=-1$ for the integral class
$[\Sigma]\in H_2(X_+;\Z)$. If $c_1(L_+)=m\,g$ for a generator $g$ of $H^2(X_+;\Z)\cong\Z$, then
$m\langle g,[\Sigma]\rangle=-1$, so $m=\pm1$.
\end{proof}

This uses only items (2)--(6) above. It is consistent with \cite[Lemma 4.1]{LLZ26}: for
instance $L_-|_{\mathcal F}\cong\cO(-1,-1)$, so $c_1(L_-)^2\cdot[\mathcal F]=2=2\,c_1(L_-)^3$.

\begin{remark}[the complex $S^2\times S^4$ of Guo--Fang--Lu]
For the complex $S^2\times S^4$ of \cite{GFL26} we did not find a holomorphic line bundle
with primitive Chern class. The two components $D_0,D_1$ of its cusp fibre satisfy
$[D_0]=2x=-[D_1]$ in $H^2(S^2\times S^4;\Z)=\Z x$ (end of the proof of
\cite[Prop.~2.2]{GFL26}). Applied to $\cO(D_0)$, the construction below would give a
lens-space factor $L(2,1)$ in place of $S^3$. Theorem~C therefore uses \cite{LLZ26}.
\end{remark}

\subsection{Proof of Theorem C}\label{sec:proofC}
\begin{proof}[Proof of Theorem C]
Assume (I3), and fix $\tau\in\C\setminus\R$, for instance $\tau=i$.

\emph{$S^4\times S^3\times S^{2q+1}$, $q\ge0$.}
\begin{itemize}
\item Let $B=X_+\times\PP^q$, a compact connected complex manifold of dimension $3+q$.
\item Let $L_1=\mathrm{pr}_1^*L_+$ and $L_2=\mathrm{pr}_2^*\cO_{\PP^q}(-1)$. For $q=0$, $L_2$ is
  the trivial bundle.
\item By Proposition~\ref{prop:CE}, $Z_q=(L_1^\times\times_BL_2^\times)/\C$ is a compact
  complex manifold of dimension $q+4$, and a holomorphic $E_\tau$-bundle over $B$.
\item Also by Proposition~\ref{prop:CE}, it is diffeomorphic to
  \[
  S(L_1)\times_BS(L_2)\ \cong\ S(L_+)\times S(\cO_{\PP^q}(-1))\ \cong\ (S^3\times S^4)\times S^{2q+1}.
  \]
  The first step is Lemma~\ref{lem:productcircle}. The second uses
  Lemma~\ref{lem:S3S4}, which applies by (I3) and Lemma~\ref{lem:primitive}, together
  with (C4).
\item As a check, the real dimensions agree: $3+4+(2q+1)=2(q+4)$.
\end{itemize}

\emph{$(S^4\times S^3)^2$.} Let $B=X_+\times X_+$ and $L_j=\mathrm{pr}_j^*L_+$.
Proposition~\ref{prop:CE} and Lemmas~\ref{lem:productcircle} and~\ref{lem:S3S4} give a
compact complex manifold of dimension $7$, diffeomorphic to
$S(L_+)\times S(L_+)\cong(S^3\times S^4)^2$.
\end{proof}

\begin{remark}\label{rem:Zq}
\begin{enumerate}
\item \emph{The case $q=0$.} By Remark~\ref{rem:onebundle}, for $\tau=2\pi i/\log2$ one has
  $Z_0\cong L_+^\times/2^{\Z}$, the
  quotient of $L_+$ minus its zero section by $v\mapsto2v$ (for $\tau=i$, replace $2$ by
  $e^{2\pi}$). The map
  $v\mapsto(v/|v|,\ \log|v|\bmod\log2)$ identifies it with
  $S(L_+)\times S^1\cong S^4\times S^3\times S^1$.
\item \emph{Relation to Calabi--Eckmann.} With $X_+$ replaced by $\PP^p$ and $L_+$ by
  $\cO(-1)$, the construction is exactly the Calabi--Eckmann manifold
  $S^{2p+1}\times S^{2q+1}$. The only new ingredient is the base $X_+$ with its line bundle
  $L_+$.
\item \emph{Not Kähler.} $H^2(S^4\times S^3\times S^{2q+1};\R)=0$ for every $q\ge0$, so
  $Z_q$ is not Kähler.
\item \emph{Canonical bundle.} $Z_q\to X_+\times\PP^q$ is a holomorphic principal
  $E_\tau$-bundle. Its vertical tangent bundle is trivial, so
  $K_{Z_q}\cong\pi^*K_{X_+\times\PP^q}$.
\item \emph{Why this does not give $S^4\times T^2$.} Here every $S^4$-factor comes from
  $X_+\cong S^2\times S^4$, so it arrives paired with the $S^3$ of a Hopf circle over the
  $S^2$-factor. To obtain $S^4\times T^2$ in this way, one would need a compact complex
  surface carrying a principal $T^2$-bundle with total space $S^4\times T^2$. We know of
  none, and the obvious base $S^4$ is not complex. The Main Theorem is therefore
  genuinely needed for $S^4\times T^2$ and $S^4\times T^4$ below.
\end{enumerate}
\end{remark}

\subsection{Products of spheres of dimension at most eight}\label{sec:dim8}
\begin{lemma}\label{lem:noacs}
$S^4$, $S^8$ and $S^4\times S^4$ admit no almost complex structure.
\end{lemma}
\begin{proof}
Let $M$ be one of these manifolds, of dimension $2m\in\{4,8\}$.
\begin{itemize}
\item \emph{Pontryagin classes vanish.} $TS^n\oplus\underline\R\cong\underline\R^{n+1}$, so
  $M$ is stably parallelisable and all its Pontryagin classes vanish.
\item \emph{Chern classes.} Suppose $J$ is an almost complex structure with Chern classes
  $c_i=c_i(TM,J)$. Since $TM\otimes\C\cong(TM,J)\oplus\overline{(TM,J)}$, we have
  $\sum_k(-1)^kp_k=c(TM,J)\,c(\overline{TM,J})$. In low degrees this gives
  \[p_1=c_1^2-2c_2,\qquad p_2=c_2^2-2c_1c_3+2c_4.\]
  Moreover $c_m=e(TM)$ for the orientation of $J$, and $\langle e(TM),[M]\rangle=\chi(M)$,
  which is $2$, $2$ and $4$ respectively.
\item \emph{$M=S^4$.} $c_1=0$ since $H^2=0$. Then $0=p_1=-2c_2$ forces $c_2=0$, since
  $H^4(S^4;\Z)\cong\Z$. But $c_2=e\neq0$.
\item \emph{$M=S^8$ or $S^4\times S^4$.} $c_1=c_3=0$ since $H^2=H^6=0$. Then $0=p_1=-2c_2$
  forces $c_2=0$, since $H^4$ is torsion-free. Then $0=p_2=2c_4=2e$. But $2e\neq0$ in
  $H^8(M;\Z)\cong\Z$.
\end{itemize}
\end{proof}

This recovers the case $S^4$ of Ehresmann and Hopf, the case $S^8$ of Borel--Serre
\cite{BorelSerre53}, and the case $S^4\times S^4$ of Sutherland \cite{Sutherland65} and
Datta--Subramanian \cite{DattaSubramanian90}.

\begin{proof}[Proof of Corollary D]
\emph{The list.} The products of spheres of dimension $2,4,6,8$ correspond to the
partitions of these numbers. There are $2$, $5$, $11$ and $22$ of them, $40$ in all, and
they are exactly the rows of Table~\ref{tab:spheres}.

\emph{Not almost complex.} By Lemma~\ref{lem:noacs}, $S^4$, $S^8$ and $S^4\times S^4$ admit
no almost complex structure.

\emph{Complex.} Each of the other $37$ products is diffeomorphic to the complex manifold
in the second column of the table. A product of complex manifolds is complex. The
notation is:
\begin{itemize}
\item $\mathcal H_k=(\C^k\setminus0)/2^{\Z}\cong S^{2k-1}\times S^1$ is a Hopf manifold;
\item $\mathrm{CE}_{p,q}\cong S^{2p+1}\times S^{2q+1}$ is a Calabi--Eckmann manifold
  \cite{CalabiEckmann53};
\item $E$ is an elliptic curve, so $E\cong T^2$;
\item $X_{S^2\times S^4}$ is either complex $S^2\times S^4$ of (I2);
\item $Y\cong S^4\times T^2$ is the threefold of the Main Theorem;
\item $Z_0\cong S^4\times S^3\times S^1$ is from Theorem~C.
\end{itemize}
The third column records which input each entry uses. Hence, given (I1)--(I3) and the
Main Theorem, every product of spheres of dimension at most eight other than $S^4$, $S^8$
and $S^4\times S^4$ is complex, and these three are not almost complex.
\end{proof}

\begin{table}[ht]
\centering
\small
\renewcommand{\arraystretch}{1.08}
\begin{tabular}{@{}llll@{}}
\toprule
dim & product & complex manifold & input\\
\midrule
2 & $S^2$ & $\PP^1$ & classical\\
  & $T^2$ & $E$ & classical\\
\midrule
4 & $S^4$ & none (Lemma~\ref{lem:noacs}) & ---\\
  & $S^3\times S^1$ & $\mathcal H_2$ & classical\\
  & $S^2\times S^2$, $S^2\times T^2$, $T^4$ & $(\PP^1)^2$, $\PP^1\times E$, $E^2$ & classical\\
\midrule
6 & $S^5\times S^1$, $S^3\times S^3$ & $\mathcal H_3$, $\mathrm{CE}_{1,1}$ & classical\\
  & $S^3\times S^2\times S^1$, $S^3\times T^3$ & $\mathcal H_2\times\PP^1$, $\mathcal H_2\times E$ & classical\\
  & $(S^2)^3$, $(S^2)^2\times T^2$, $S^2\times T^4$, $T^6$ & products of $\PP^1$ and $E$ & classical\\
  & $S^6$ & $X_{\rm sph}$ & (I1)\\
  & $S^4\times S^2$ & $X_{S^2\times S^4}$ & (I2)\\
  & $S^4\times T^2$ & $Y$ & Main Theorem\\
\midrule
8 & $S^8$, $S^4\times S^4$ & none (Lemma~\ref{lem:noacs}) & ---\\
  & $S^7\times S^1$, $S^5\times S^3$ & $\mathcal H_4$, $\mathrm{CE}_{2,1}$ & classical\\
  & $S^5\times S^2\times S^1$, $S^5\times T^3$ & $\mathcal H_3\times\PP^1$, $\mathcal H_3\times E$ & classical\\
  & $S^3\times S^3\times S^2$, $S^3\times S^3\times T^2$ & $\mathrm{CE}_{1,1}\times\PP^1$, $\mathrm{CE}_{1,1}\times E$ & classical\\
  & $S^3\times S^2\times S^2\times S^1$, $S^3\times S^2\times T^3$, $S^3\times T^5$ & $\mathcal H_2\times(\PP^1)^2$, $\mathcal H_2\times\PP^1\times E$, $\mathcal H_2\times E^2$ & classical\\
  & $(S^2)^4$, $(S^2)^3\times T^2$, $(S^2)^2\times T^4$, $S^2\times T^6$, $T^8$ & products of $\PP^1$ and $E$ & classical\\
  & $S^6\times S^2$, $S^6\times T^2$ & $X_{\rm sph}\times\PP^1$, $X_{\rm sph}\times E$ & (I1)\\
  & $S^4\times S^2\times S^2$ & $X_{S^2\times S^4}\times\PP^1$ & (I2)\\
  & $S^4\times S^2\times T^2$ & $X_{S^2\times S^4}\times E$, or $Y\times\PP^1$ & (I2) or Main Thm.\\
  & $S^4\times S^3\times S^1$ & $Z_0$ & (I3), Theorem~C\\
  & $S^4\times T^4$ & $Y\times E$ & Main Theorem\\
\bottomrule
\end{tabular}
\medskip
\caption{The $40$ products of spheres of dimension at most eight.}\label{tab:spheres}
\end{table}

\begin{remark}[what was open]\label{rem:wasopen}
Before 2026, the almost complex products of spheres of dimension at most eight with no
known complex structure were, to our knowledge, the nine rows of Table~\ref{tab:spheres}
not marked ``classical'' or ``---''. They are almost complex: $S^6$ and $S^2\times S^4$
classically, and the others as products of almost complex manifolds or by Staples'
theorem \cite{Staples67}.
\begin{itemize}
\item The constructions of \cite{AC26,Engel26,GFL26,LLZ26} settle $S^6$,
  $S^2\times S^4$ and their products with $\PP^1$ and $E$.
\item The Main Theorem settles $S^4\times T^2$ and $S^4\times T^4$.
\item Theorem~C settles $S^4\times S^3\times S^1$.
\end{itemize}
Without the Main Theorem, $S^4\times T^2$ would be the only open case in dimension $\le6$,
and $S^4\times T^4$ the only one in dimension $8$.
\end{remark}

\subsection{Higher dimensions}\label{sec:higherdim}
For products of even-dimensional spheres, Sutherland's theorem
\cite[Thm.~3.1]{Sutherland65} (as stated in \cite{LLZ26}; \cite{DattaSubramanian90} treats
two factors) says that the almost complex ones are exactly the products of copies of
$S^2$, $S^6$ and $S^2\times S^4$. Given (I1) and (I2), all of these are complex
\cite[Cor.~1.3]{LLZ26}. Products with an odd-dimensional factor are parallelisable
\cite{Staples67}, hence almost complex.

The known complex blocks are:
\begin{itemize}
\item $S^2$, $S^6$, $S^2\times S^4$;
\item $S^{2p+1}\times S^{2q+1}$ (Calabi--Eckmann, which includes Hopf manifolds and $T^2$);
\item $S^4\times S^3\times S^{2q+1}$ and $(S^4\times S^3)^2$ (Theorem~C);
\item $S^4\times T^2$ (Main Theorem).
\end{itemize}
No known block contains an $S^{2n}$ with $n\ge4$. Each $S^4$-factor must be absorbed in one
of four ways:
\begin{itemize}
\item by an $S^2$-factor;
\item by two $S^1$-factors;
\item by an $S^3$-factor and one further odd sphere;
\item together with a second $S^4$-factor, by two $S^3$-factors.
\end{itemize}

An exhaustive search over these blocks (Appendix~\ref{app:spheres}) finds the following
almost complex products with no complex structure obtainable in this way:
\begin{itemize}
\item dimension $10$: $S^8\times T^2$, $S^5\times S^4\times S^1$, $S^4\times S^4\times T^2$;
\item dimension $12$: $S^{10}\times T^2$, $S^8\times S^3\times S^1$, $S^8\times S^2\times T^2$,
  $S^8\times T^4$, $S^7\times S^4\times S^1$, $S^4\times S^4\times S^3\times S^1$.
\end{itemize}
The simplest of these, $S^{2n}\times T^2$ for $n\ge4$, are the higher-dimensional analogues
of Calabi's question.

%======================================================================
\section{Relation to the literature}\label{sec:literature}
%======================================================================

To the best of our knowledge, no complex structure on $S^4\times T^2$ has previously
been constructed or claimed. We summarise the work most closely related to the
present construction.

\subsection*{Calabi's question}
\begin{itemize}
\item Calabi posed the question for every closed orientable surface $V^2$
  \cite[p.~408]{Calabi58}. The torus case is the genus-one instance.
\item MathOverflow question 89798 \cite{MO89798} (2012) extends the question about
  $S^2\times S^4$ explicitly to $\Sigma\times S^4$ for all compact Riemann surfaces
  $\Sigma$. It was recorded there as open in 2015.
\item Euh and Sekigawa prove non-integrability only for \emph{orthogonal} almost
  complex structures with respect to special metrics. These are the round product
  \cite{EuhSekigawa13}, and warped products of a closed surface of nonnegative
  curvature (hence including flat $T^2$) with a round $S^4$ \cite{EuhSekigawa16}.
\item The 2026 constructions on $S^2\times S^4$ treat $V=S^2$ only \cite{GFL26,LLZ26}.
\end{itemize}

\subsection*{The 2026 constructions}
\begin{itemize}
\item The constructions on $S^6$ \cite{AC26,Engel26} and on $S^2\times S^4$
  \cite{GFL26,LLZ26}, and the subsequent constructions of which we are aware (complex
  structures on $S^6$, on rational homology spheres, and on torus-fibred threefolds
  with finite fundamental group), all have finite fundamental group.
\item None of them treats $S^4\times T^2$. None uses two cusps of opposite chirality
  or a multiplicity-one logarithmic transform.
\item Engel remarks that it is unclear which $6$-manifolds the torus-fibration method
  produces \cite[Rem.~8.5]{Engel26}.
\item These works depend on one another. \cite{GFL26} reuse the analytic framework of
  \cite{AC26}, and \cite{LLZ26} explicitly assume the validity of the construction of
  \cite{AC26}; see \S\ref{sec:inputs}. We cite them for their stated results and have
  not independently verified their proofs. The Main Theorem and Theorem~B do not use
  any of them. Theorem~C and Corollary~D do, as recorded in \S\ref{sec:inputs}.
\end{itemize}

\subsection*{Classical families}
We know of no classical family of non-K\"ahler manifolds that contains a manifold
homotopy equivalent to $S^4\times T^2$. Recall that $\pi_1(S^4\times T^2)=\Z^2$,
$\pi_2=0$, and the Betti numbers are $1,2,1,0,1,2,1$. For each family we name the
invariant that rules it out.
\begin{itemize}
\item \emph{Calabi--Eckmann, Hopf and Brieskorn--Van de Ven manifolds.} In real dimension
  six these are products of two odd-dimensional (homotopy) spheres, so $b_2=0$.
\item \emph{LVM manifolds.} An LVM manifold $N$ is either a projective space or the
  quotient of a link $X_A$ by the free diagonal circle action
  \cite[\S12, discussion before Thm.~12.2]{BosioMeersseman06}.
  \begin{itemize}
  \item By \cite[Lemma~0.9]{BosioMeersseman06}, $X_A$ is equivariantly diffeomorphic
    to $X_B\times(S^1)^k$, with $X_B$ a $2$-connected link and $k$ the rank of
    $\pi_1(X_A)$.
  \item If $k=0$, the exact sequence of $S^1\to X_A\to N$ gives $\pi_1(N)=0$.
  \item If $k\ge1$, then $N\cong X_B\times T^{k-1}$ (loc.\ cit., \S12). So
    $\pi_1(N)=\Z^2$ forces $k=3$ and $N\cong X_B\times T^2$, with $X_B$ a closed
    $2$-connected $4$-manifold.
  \item By Poincar\'e duality such an $X_B$ has $\chi(X_B)=2$. But every link has
    Euler characteristic $0$ \cite[Lemma~0.10]{BosioMeersseman06}.
  \end{itemize}
  So no LVM manifold has fundamental group $\Z^2$ in complex dimension three.
\item \emph{Maximal torus actions.} Compact complex manifolds with maximal torus
  actions are classified \cite{Ishida19}. The standard $T^4$-action on $S^4\times T^2$
  does not occur there, since two of its characteristic submanifolds meet in a
  disconnected set. This excludes only that action. We have not determined whether
  $S^4\times T^2$ carries another maximal torus action compatible with a complex
  structure. Our $T^2$-action is not maximal.
\item \emph{Twistor spaces.} A twistor space $Z\to M^4$ has
  $H^*(Z;\Q)\cong H^*(M;\Q)\otimes H^*(S^2;\Q)$ by Leray--Hirsch. So
  $b_3(Z)=b_3(M)+b_1(M)$, and $b_1(Z)=2$ forces $b_3(Z)\ge2$, whereas
  $b_3(S^4\times T^2)=0$.
\item \emph{Vaisman manifolds} have odd $b_1$ \cite{Vaisman82}.
\item \emph{Kato manifolds}, i.e.\ compact complex manifolds with a global spherical
  shell, have $\pi_1\cong\Z$ \cite{Kato77}; see \cite[\S1]{IOP19}.
\item \emph{Compact homogeneous complex manifolds.} Such an $X$ fibres over a
  rational homogeneous manifold $Q$ with a compact complex parallelisable fibre
  $F=\Gamma\backslash L$ \cite{Tits62,Wang54}; $F$ is connected since $Q$ is simply
  connected. Suppose $X\simeq S^4\times T^2$.
  \begin{itemize}
  \item If $Q$ is a point, then the universal cover $\widetilde X\simeq S^4$ is
    homotopy equivalent to the universal cover of the complex Lie group $L$. That cover
    is either contractible or has $H^3\ne0$, a contradiction.
  \item Otherwise $\dim_\C F\le2$. A complex Lie group of dimension $\le2$ with a
    cocompact lattice is abelian, so $F$ is a torus $T^{2d}$. The exact sequence
    $0=\pi_2(X)\to\pi_2(Q)\to\pi_1(F)\to\pi_1(X)\to0$ gives $2d=b_2(Q)+2$, with
    $b_2(Q)\ge1$ and $\dim_\C Q=3-d$. Hence $d=2$ and $Q=\PP^1$, but then
    $b_2(Q)=1\ne2$.
  \end{itemize}
\item $S^4\times T^2$ is not symplectic, since $H^2$ is generated by a class $a$ with
  $a^2=0$. So any complex structure on it is non-K\"ahler.
\end{itemize}

\subsection*{Products of spheres}
\begin{itemize}
\item Only $S^2$ and $S^6$ are almost complex among spheres \cite{BorelSerre53}.
\item Among products of even-dimensional spheres, the almost complex ones are the
  products of copies of $S^2$, $S^6$ and $S^2\times S^4$ \cite[Thm.~3.1]{Sutherland65}. We
  rely on the statement in \cite{LLZ26}; for two factors see also
  \cite{DattaSubramanian90}.
\item A product of spheres with an odd-dimensional factor is parallelisable
  \cite{Staples67}.
\item Products of two odd-dimensional spheres are complex \cite{CalabiEckmann53}.
  Calabi--Eckmann type constructions on torus bundles over complex bases are standard
  \cite{Morimoto63,Hofer93,GoldsteinProkushkin04,GGP08}.
\item Given the complex $S^6$, the complex $S^2\times S^4$ of \cite{LLZ26} makes every
  almost complex product of even-dimensional spheres complex \cite[Cor.~1.3]{LLZ26}.
\end{itemize}

\subsection*{The closest antecedent}
Di Scala, Kasuya and Zuddas \cite{DKZ17} make the positive half of the
Matsumoto--Fukaya fibration holomorphic on $\R^4\subset S^4$. This gives non-K\"ahler
complex structures on $\R^4$, and they show that these surfaces admit no smooth
complex compactification. The present construction realises \emph{both} nodes. It
does so after fibring by $T^2$, using $2$-tori that are very generally non-algebraic, and in compact complex
dimension three.

%%%%%%%%%%%%%%%%%%%%%%%%%%%%%%%%%%%%%%%%%%%%%%%%%
\section*{Acknowledgements}
EH is supported by S\~ao Paulo Research Foundation (FAPESP) grant 2024/18994-7.
TSG is supported by the 2024 Max Planck--Humboldt Research Award bestowed on Geordie Williamson and Catharina Stroppel by the Max Planck Society and the Alexander von Humboldt Foundation.
The symbolic and numerical verifications were carried out in Python (\texttt{sympy}, \texttt{numpy}, \texttt{scipy}); the computations and the first draft of this note were prepared with the assistance of Claude (Anthropic).

%======================================================================
\appendix
\section{Symbolic verification}\label{app:checks}
%======================================================================

The accompanying \texttt{sympy} script \texttt{verify\_construction.py}, in the
\texttt{code/} folder of the source files, checks the following items in exact symbolic
arithmetic.
\begin{enumerate}
\item The determinant of Lemma~\ref{lem:det-family}; it is independent of $\arg t$.
\item The monodromy $e_4\mapsto e_4+e_1$.
\item The conjugation formulas of Lemma~\ref{lem:atinfty}.
\item $\tilde\sigma(s+1)-\tilde\sigma(s)=e_4(s)$ (Lemma~\ref{lem:Phi}).
\item The extension of the $T^2$-action across $\infty$ (Proposition~\ref{prop:torus}).
\item The Jacobian determinants $\partial(z,w)/\partial(z'',w')=s^{-2}$ on
  fibres, and $\partial(z,w,t)/\partial(z'',w',s)=-s^{-4}$.
\item The determinant of Remark~\ref{rem:deform}.
\item The toric chart identities of Lemma~\ref{lem:toric}(1),(5),(6), including the
  dual-basis identities $A^Tm^{(n+1)}_i=m^{(n)}_i$.
\item $\omega(\eb_2,\eb_3)=1-|t|^4$ near $0$ and $|s|^4-1$ near $\infty$.
\item Lemma~\ref{lem:localdet} for arbitrary symbolic data.
\item The exact heights $P_0$ and $P_\infty$ of Remark~\ref{rem:Pexact}, computed
  from the filling data, and the resulting bounds at $\varepsilon=10^{-2}$ and
  $\varepsilon=1/16$.
\item The periods of Lemma~\ref{lem:periods}, the absence of integral relations among
  them, and the product decomposition of $X_1$ in Proposition~\ref{prop:NS}.
\item The derivative of the period ratio $p_{13}/p_{12}$ used in \S\ref{sec:props}.
\item The critical point, the maximum and the counterexample in the description of
  the admissible region \eqref{eq:admissible}.
\end{enumerate}

\subsection{Enumeration of products of spheres}\label{app:spheres}
The accompanying script \texttt{sphere\_products.py}, also in \texttt{code/}, does the
following:
\begin{enumerate}
\item it enumerates the products of spheres of dimensions $2,4,\dots,12$, i.e.\ the
  partitions of these numbers;
\item it decides almost complexity by Sutherland's theorem (all factors even) or
  Staples' theorem (some factor odd);
\item it searches all decompositions of each almost complex product into the known
  complex blocks of \S\ref{sec:higherdim};
\item it records the inclusion-minimal sets of inputs (I1), (I2), (I3) and the Main
  Theorem that each product needs.
\end{enumerate}
It asserts the classification of Corollary~D in dimension $8$. In particular, it checks
that $S^4\times S^3\times S^1$ needs (I3), that $S^4\times T^4$ needs the Main Theorem, and
that no almost complex product of dimension $8$ is left over. It also produces the
lists of \S\ref{sec:higherdim}.

%%%%%%%%%%%%%%%%%%%%%%%%%%%%%%%%%%%%%%%%%%%%%%%%%
%\newpage
\addcontentsline{toc}{section}{References}
\bibliographystyle{utphys}
\bibliography{references}{}

\end{document}
